% Options for packages loaded elsewhere
% Options for packages loaded elsewhere
\PassOptionsToPackage{unicode}{hyperref}
\PassOptionsToPackage{hyphens}{url}
\PassOptionsToPackage{dvipsnames,svgnames,x11names}{xcolor}
%
\documentclass[
  english,
  letterpaper,
  DIV=11,
  numbers=noendperiod]{scrreprt}
\usepackage{xcolor}
\usepackage{amsmath,amssymb}
\setcounter{secnumdepth}{5}
\usepackage{iftex}
\ifPDFTeX
  \usepackage[T1]{fontenc}
  \usepackage[utf8]{inputenc}
  \usepackage{textcomp} % provide euro and other symbols
\else % if luatex or xetex
  \usepackage{unicode-math} % this also loads fontspec
  \defaultfontfeatures{Scale=MatchLowercase}
  \defaultfontfeatures[\rmfamily]{Ligatures=TeX,Scale=1}
\fi
\usepackage{lmodern}
\ifPDFTeX\else
  % xetex/luatex font selection
\fi
% Use upquote if available, for straight quotes in verbatim environments
\IfFileExists{upquote.sty}{\usepackage{upquote}}{}
\IfFileExists{microtype.sty}{% use microtype if available
  \usepackage[]{microtype}
  \UseMicrotypeSet[protrusion]{basicmath} % disable protrusion for tt fonts
}{}
\makeatletter
\@ifundefined{KOMAClassName}{% if non-KOMA class
  \IfFileExists{parskip.sty}{%
    \usepackage{parskip}
  }{% else
    \setlength{\parindent}{0pt}
    \setlength{\parskip}{6pt plus 2pt minus 1pt}}
}{% if KOMA class
  \KOMAoptions{parskip=half}}
\makeatother
% Make \paragraph and \subparagraph free-standing
\makeatletter
\ifx\paragraph\undefined\else
  \let\oldparagraph\paragraph
  \renewcommand{\paragraph}{
    \@ifstar
      \xxxParagraphStar
      \xxxParagraphNoStar
  }
  \newcommand{\xxxParagraphStar}[1]{\oldparagraph*{#1}\mbox{}}
  \newcommand{\xxxParagraphNoStar}[1]{\oldparagraph{#1}\mbox{}}
\fi
\ifx\subparagraph\undefined\else
  \let\oldsubparagraph\subparagraph
  \renewcommand{\subparagraph}{
    \@ifstar
      \xxxSubParagraphStar
      \xxxSubParagraphNoStar
  }
  \newcommand{\xxxSubParagraphStar}[1]{\oldsubparagraph*{#1}\mbox{}}
  \newcommand{\xxxSubParagraphNoStar}[1]{\oldsubparagraph{#1}\mbox{}}
\fi
\makeatother

\usepackage{color}
\usepackage{fancyvrb}
\newcommand{\VerbBar}{|}
\newcommand{\VERB}{\Verb[commandchars=\\\{\}]}
\DefineVerbatimEnvironment{Highlighting}{Verbatim}{commandchars=\\\{\}}
% Add ',fontsize=\small' for more characters per line
\usepackage{framed}
\definecolor{shadecolor}{RGB}{241,243,245}
\newenvironment{Shaded}{\begin{snugshade}}{\end{snugshade}}
\newcommand{\AlertTok}[1]{\textcolor[rgb]{0.68,0.00,0.00}{#1}}
\newcommand{\AnnotationTok}[1]{\textcolor[rgb]{0.37,0.37,0.37}{#1}}
\newcommand{\AttributeTok}[1]{\textcolor[rgb]{0.40,0.45,0.13}{#1}}
\newcommand{\BaseNTok}[1]{\textcolor[rgb]{0.68,0.00,0.00}{#1}}
\newcommand{\BuiltInTok}[1]{\textcolor[rgb]{0.00,0.23,0.31}{#1}}
\newcommand{\CharTok}[1]{\textcolor[rgb]{0.13,0.47,0.30}{#1}}
\newcommand{\CommentTok}[1]{\textcolor[rgb]{0.37,0.37,0.37}{#1}}
\newcommand{\CommentVarTok}[1]{\textcolor[rgb]{0.37,0.37,0.37}{\textit{#1}}}
\newcommand{\ConstantTok}[1]{\textcolor[rgb]{0.56,0.35,0.01}{#1}}
\newcommand{\ControlFlowTok}[1]{\textcolor[rgb]{0.00,0.23,0.31}{\textbf{#1}}}
\newcommand{\DataTypeTok}[1]{\textcolor[rgb]{0.68,0.00,0.00}{#1}}
\newcommand{\DecValTok}[1]{\textcolor[rgb]{0.68,0.00,0.00}{#1}}
\newcommand{\DocumentationTok}[1]{\textcolor[rgb]{0.37,0.37,0.37}{\textit{#1}}}
\newcommand{\ErrorTok}[1]{\textcolor[rgb]{0.68,0.00,0.00}{#1}}
\newcommand{\ExtensionTok}[1]{\textcolor[rgb]{0.00,0.23,0.31}{#1}}
\newcommand{\FloatTok}[1]{\textcolor[rgb]{0.68,0.00,0.00}{#1}}
\newcommand{\FunctionTok}[1]{\textcolor[rgb]{0.28,0.35,0.67}{#1}}
\newcommand{\ImportTok}[1]{\textcolor[rgb]{0.00,0.46,0.62}{#1}}
\newcommand{\InformationTok}[1]{\textcolor[rgb]{0.37,0.37,0.37}{#1}}
\newcommand{\KeywordTok}[1]{\textcolor[rgb]{0.00,0.23,0.31}{\textbf{#1}}}
\newcommand{\NormalTok}[1]{\textcolor[rgb]{0.00,0.23,0.31}{#1}}
\newcommand{\OperatorTok}[1]{\textcolor[rgb]{0.37,0.37,0.37}{#1}}
\newcommand{\OtherTok}[1]{\textcolor[rgb]{0.00,0.23,0.31}{#1}}
\newcommand{\PreprocessorTok}[1]{\textcolor[rgb]{0.68,0.00,0.00}{#1}}
\newcommand{\RegionMarkerTok}[1]{\textcolor[rgb]{0.00,0.23,0.31}{#1}}
\newcommand{\SpecialCharTok}[1]{\textcolor[rgb]{0.37,0.37,0.37}{#1}}
\newcommand{\SpecialStringTok}[1]{\textcolor[rgb]{0.13,0.47,0.30}{#1}}
\newcommand{\StringTok}[1]{\textcolor[rgb]{0.13,0.47,0.30}{#1}}
\newcommand{\VariableTok}[1]{\textcolor[rgb]{0.07,0.07,0.07}{#1}}
\newcommand{\VerbatimStringTok}[1]{\textcolor[rgb]{0.13,0.47,0.30}{#1}}
\newcommand{\WarningTok}[1]{\textcolor[rgb]{0.37,0.37,0.37}{\textit{#1}}}

\usepackage{longtable,booktabs,array}
\usepackage{calc} % for calculating minipage widths
% Correct order of tables after \paragraph or \subparagraph
\usepackage{etoolbox}
\makeatletter
\patchcmd\longtable{\par}{\if@noskipsec\mbox{}\fi\par}{}{}
\makeatother
% Allow footnotes in longtable head/foot
\IfFileExists{footnotehyper.sty}{\usepackage{footnotehyper}}{\usepackage{footnote}}
\makesavenoteenv{longtable}
\usepackage{graphicx}
\makeatletter
\newsavebox\pandoc@box
\newcommand*\pandocbounded[1]{% scales image to fit in text height/width
  \sbox\pandoc@box{#1}%
  \Gscale@div\@tempa{\textheight}{\dimexpr\ht\pandoc@box+\dp\pandoc@box\relax}%
  \Gscale@div\@tempb{\linewidth}{\wd\pandoc@box}%
  \ifdim\@tempb\p@<\@tempa\p@\let\@tempa\@tempb\fi% select the smaller of both
  \ifdim\@tempa\p@<\p@\scalebox{\@tempa}{\usebox\pandoc@box}%
  \else\usebox{\pandoc@box}%
  \fi%
}
% Set default figure placement to htbp
\def\fps@figure{htbp}
\makeatother



\ifLuaTeX
\usepackage[bidi=basic]{babel}
\else
\usepackage[bidi=default]{babel}
\fi
% get rid of language-specific shorthands (see #6817):
\let\LanguageShortHands\languageshorthands
\def\languageshorthands#1{}
\ifLuaTeX
  \usepackage[english]{selnolig} % disable illegal ligatures
\fi


\setlength{\emergencystretch}{3em} % prevent overfull lines

\providecommand{\tightlist}{%
  \setlength{\itemsep}{0pt}\setlength{\parskip}{0pt}}



 


\usepackage{booktabs}
\usepackage{longtable}
\usepackage{array}
\usepackage{multirow}
\usepackage{wrapfig}
\usepackage{float}
\usepackage{colortbl}
\usepackage{pdflscape}
\usepackage{tabu}
\usepackage{threeparttable}
\usepackage{threeparttablex}
\usepackage[normalem]{ulem}
\usepackage{makecell}
\usepackage{xcolor}
\KOMAoption{captions}{tableheading}
\makeatletter
\@ifpackageloaded{bookmark}{}{\usepackage{bookmark}}
\makeatother
\makeatletter
\@ifpackageloaded{caption}{}{\usepackage{caption}}
\AtBeginDocument{%
\ifdefined\contentsname
  \renewcommand*\contentsname{Table of contents}
\else
  \newcommand\contentsname{Table of contents}
\fi
\ifdefined\listfigurename
  \renewcommand*\listfigurename{List of Figures}
\else
  \newcommand\listfigurename{List of Figures}
\fi
\ifdefined\listtablename
  \renewcommand*\listtablename{List of Tables}
\else
  \newcommand\listtablename{List of Tables}
\fi
\ifdefined\figurename
  \renewcommand*\figurename{Figure}
\else
  \newcommand\figurename{Figure}
\fi
\ifdefined\tablename
  \renewcommand*\tablename{Table}
\else
  \newcommand\tablename{Table}
\fi
}
\@ifpackageloaded{float}{}{\usepackage{float}}
\floatstyle{ruled}
\@ifundefined{c@chapter}{\newfloat{codelisting}{h}{lop}}{\newfloat{codelisting}{h}{lop}[chapter]}
\floatname{codelisting}{Listing}
\newcommand*\listoflistings{\listof{codelisting}{List of Listings}}
\makeatother
\makeatletter
\makeatother
\makeatletter
\@ifpackageloaded{caption}{}{\usepackage{caption}}
\@ifpackageloaded{subcaption}{}{\usepackage{subcaption}}
\makeatother
\usepackage{bookmark}
\IfFileExists{xurl.sty}{\usepackage{xurl}}{} % add URL line breaks if available
\urlstyle{same}
\hypersetup{
  pdftitle={Country Analysis Meeting (CAM) 2025 - Analysis Guide},
  pdflang={en},
  colorlinks=true,
  linkcolor={blue},
  filecolor={Maroon},
  citecolor={Blue},
  urlcolor={Blue},
  pdfcreator={LaTeX via pandoc}}


\title{Country Analysis Meeting (CAM) 2025 - Analysis Guide}
\author{}
\date{Jun 15, 2025}
\begin{document}
\maketitle

\renewcommand*\contentsname{Table of contents}
{
\hypersetup{linkcolor=}
\setcounter{tocdepth}{2}
\tableofcontents
}

\bookmarksetup{startatroot}

\chapter*{Welcome}\label{welcome}
\addcontentsline{toc}{chapter}{Welcome}

\markboth{Welcome}{Welcome}

\section*{Overview of the Countdown to
2030}\label{overview-of-the-countdown-to-2030}
\addcontentsline{toc}{section}{Overview of the Countdown to 2030}

\markright{Overview of the Countdown to 2030}

The Countdown to 2030 for Women's, Children's and Adolescents' Health
(The Countdown) initiative is a global collaboration involving academics
from national, regional, and international institutions, UN agencies,
the World Bank, and civil society organizations. The initiative tracks
progress in Reproductive, Maternal, Newborn, Child, and Adolescent
Health and Nutrition (RMNCAH+N), fostering advocacy and accountability
through rigorous data analysis.

\textbf{Key objectives include:}

\begin{itemize}
\tightlist
\item
  Strengthening country-led data analysis and monitoring
\item
  Fostering innovation and evidence generation through multi-country
  collaboration
\item
  Enhancing global measurement and monitoring Improving policy and
  program communication
\end{itemize}

For more on Countdown to 2030 initiative, visit:
\href{https://www.countdown2030.org/about}{The Countdown website}

\section*{\texorpdfstring{About the \texttt{cd2030.rmncah} R
Package}{About the cd2030.rmncah R Package}}\label{about-the-cd2030.rmncah-r-package}
\addcontentsline{toc}{section}{About the \texttt{cd2030.rmncah} R
Package}

\markright{About the \texttt{cd2030.rmncah} R Package}

The cd2030.rmncah R package and Shiny App were developed to support
evidence generation and analysis of RMNCAH indicators.

Key features include:

\begin{itemize}
\item
  User-friendly interface for data management and analysis
\item
  Tools for visualization and statistical summation
\item
  Automated report generation
\end{itemize}

\section*{Acknowledgments}\label{acknowledgments}
\addcontentsline{toc}{section}{Acknowledgments}

\markright{Acknowledgments}

The Countdown acknowledges and appreciates the contribution of the
following individual in developing the Stata codes that form the
backbone of the \texttt{cd2030.rmncah} R package and Shiny App:

\part{About this Guide}

\chapter{CD2030 CAM Approach}\label{cd2030-cam-approach}

The CD2030 for Women's, Children's and Adolescents' Health, GFF, UNICEF,
WHO, WAHO and other partners are collaborating to strengthen country-led
progress and performance reviews, such as annual health sector reviews
and midterm reviews of investment cases. This guidebook is for Countdown
country analytical teams to develop a set of national and subnational
estimates for key RMNCAH-N indicators, including equity, using five-year
time series of routine data and survey results.

Much attention is paid to obtaining a clean data set with the necessary
corrections and adjustments for known biases. Service coverage and
equity, maternal and perinatal mortality, and health service utilization
and systems performance are the main subjects, with a focus on
monitoring national and sub-national targets, as well as global targets.
The figure below shows the general overview of the CAM approach

\begin{center}
\includegraphics[width=1\linewidth,height=\textheight,keepaspectratio]{pages/../images/0-overview_cam.png}
\end{center}

\section{Organization of the
Guidebook}\label{organization-of-the-guidebook}

This guidebook is organized into seven sections, each focusing on a
specific area of data analysis related to reproductive, maternal,
newborn, child and adolescent health and nutrition (RMNCAH-N). The
guidebook provides a comprehensive approach to analyzing routine health
data and survey results, with an emphasis on data quality, coverage,
equity, and health systems performance. The seven data analysis sections
in this guidebook are:

\begin{enumerate}
\def\labelenumi{\arabic{enumi}.}
\tightlist
\item
  Section 1: Health facility data quality assessment
\item
  Section 2: National~ Analyses (Coverage \& Equity)
\item
  Section 3: Subnational analysis (Coverage \& Inequality)
\item
  Section 4: Maternal mortality, stillbirths and neonatal mortality
\item
  Section5: Curative health services utilization for sick children
\item
  Section 6: Health systems progress and performance
\item
  Section 7: Planning ahead for data use
\end{enumerate}

These sections are designed to be used in a modular way, allowing
countries to select the topics that are most relevant to their context
and data availability.

Each section has:

\begin{itemize}
\tightlist
\item
  \textbf{Why/Rationale -} the scientific basis for the analysis;
\item
  \textbf{Approach}- a step-by-step guide on how to conduct the
  analysis; and the
\item
  \textbf{Implementation -} the use of the R Shiny App for data
  visualization and interpretation.
\end{itemize}

\section{Data Sources}\label{intro-data-sources}

The Countdown CAM approach uses a variety of data sources, including:

\begin{itemize}
\tightlist
\item
  \textbf{Health facility data}: Routine health data collected from
  health facilities, including service coverage, health systems
  performance, and health service utilization.
\item
  \textbf{Surveys}: Nationally representative surveys, such as the
  Demographic and Health Surveys (DHS) and Multiple Indicator Cluster
  Surveys (MICS), which provide data on health indicators, equity, and
  health service utilization.
\item
  \textbf{Administrative data}: Data collected by government agencies,
  such as vital registration systems and health management information
  systems (HMIS), which provide information on health outcomes and
  service delivery.
\item
  \textbf{Other data sources}: Other relevant data sources, such as
  census data, population estimates, and health financing data, which
  provide additional context for the analysis.
\end{itemize}

\section{Expected outputs}\label{intro-expected-outputs}

\begin{itemize}
\tightlist
\item
  \textbf{Synthesis/poster reports} (.pdf, .doc files)
\item
  \textbf{Full country report} (national/sub-national) (.pdf, .doc
  files) - to be downloaded in sections from the Shiny App and compiled
  into a full report with analysis outputs and interpretations.
\item
  \textbf{Country analytical reproducible files} (.rds files)
\item
  \textbf{Adjusted and or summarized data files} (.csv. .dta, xlsx
  files)
\end{itemize}

Country analytical reproducible files (.rds files) are the final output
of the analysis, which can be used for further analysis and
visualization. These files contain the cleaned and processed data,
country specific analysis parameters as well as the results of the
analysis, including coverage, equity, and health systems performance
indicators.

\chapter{Getting started}\label{getting-started}

In the past, Countdown to 2030 (CD2030) analyses were conducted using
Stata, relying on a series of \texttt{.do} files. These scripts have
since been translated into an R package --- \textbf{cd2030.rmncah} ---
and integrated into a user-friendly Shiny App. This application enables
health data users to load data, conduct data quality assessments (DQA),
perform analyses, and generate insightful outputs using pre-designed
report templates.

The tool facilitates in-depth inspection of data quality and supports
subnational analysis (up to the administrative level 2) for various
indicators related to data quality and RMNCAH outcomes.

\section{Software requirements}\label{sec-software-requirements}

To use this innovative digital health analytics tool, the following are
required:

\begin{itemize}
\item
  Install the latest versions of \textbf{R} and \textbf{RStudio}
\item
  Install the \textbf{cd2030.rmncah} R package latest version-noted by
  the superscript suffix in the name of the App

  Ensure access to all required
  (\hyperref[sec-data-requirements]{datasets}) and
  (\hyperref[sec-denom-analysis-setup]{parameters})
\item
  Install \textbf{RTools} (for Windows users, to support package
  compilation)
\end{itemize}

\section{Installing R and RStudio}\label{sec-instal-r}

To begin working in R, it is necessary to install both R and RStudio. R
is the underlying programming language, while RStudio provides a
user-friendly interface that simplifies the development and execution of
R scripts. Both are freely available and widely supported across
platforms.

\textbf{Note}: R and RStudio are not the same. R is the programming
language, while RStudio is an integrated development environment (IDE)
that makes working with R easier.

To ensure smooth installation and functionality, it is recommended to
install \textbf{R} first, followed by \textbf{RStudio}. This order is
important as RStudio relies on R to function properly.

\subsection{Step 1: Download and Install
R}\label{step-1-download-and-install-r}

R is distributed via The
\href{https://cran.r-project.org/}{Comprehensive R Archive Network
(CRAN)}. Select your operating system from the homepage: Windows, Mac,
or Linux.

\textbf{Windows}

\begin{enumerate}
\def\labelenumi{\arabic{enumi}.}
\tightlist
\item
  Navigate to \textbf{\emph{Download R for Windows}} and select the
  ``base'' option.
\item
  Click the first link (e.g., ``Download R x.x.x for Windows'') to
  download the installer.
\item
  Run the installer and follow the prompts. \textbf{\emph{Administrator
  privileges may be required}}.
\item
  R will be installed in your system's Program Files, with a shortcut
  added to the Start menu.
\end{enumerate}

\textbf{Mac}

\begin{enumerate}
\def\labelenumi{\arabic{enumi}.}
\tightlist
\item
  Click \textbf{\emph{Download R for Mac}} on the CRAN homepage.
\item
  Download the latest release package and run the installer.
\item
  The default installation settings are typically sufficient. You may be
  prompted to enter your system password.
\end{enumerate}

\textbf{Note}

R is not a graphical software application but a programming environment.
It is best used in conjunction with RStudio, which provides a consistent
and user-friendly interface across operating systems.

\subsection{Step 2: Download and Install
RStudio}\label{step-2-download-and-install-rstudio}

RStudio is an integrated development environment (IDE) designed for R.
It features a script editor, console, graphics viewer, and additional
tools for package management, debugging, and file organization.

Download RStudio from \href{https://posit.co/downloads/}{Posit Website}

\textbf{Do I still need to download R?} Even if you use RStudio, you'll
still need to download R to your computer. RStudio helps you use the
version of R that lives on your computer, but it doesn't come with a
version of R on its own.

After installation, launch RStudio and begin interacting with R through
its console and script windows.

\begin{center}
\includegraphics[width=1\linewidth,height=\textheight,keepaspectratio]{pages/../images/rstudio.png}
\end{center}

\section{Installing the CD2030 RMNCAH
package}\label{sec-instal-cd2030-rmncah}

After installing R and RStudio, you may proceed with installing the
Countdown2030 RMNCAH application, which is hosted on GitHub under the
repository
\href{https://github.com/aphrcwaro/cd2030.rmncah}{\texttt{cd2030.rnncah}}.
The application is implemented as an R package and supports interactive
dashboard generation via Shiny.

\subsection{Installation via R
Console}\label{installation-via-r-console}

\textbf{Stable Version}

\begin{Shaded}
\begin{Highlighting}[]
\ControlFlowTok{if}\NormalTok{ (}\SpecialCharTok{!}\FunctionTok{requireNamespace}\NormalTok{(}\StringTok{"devtools"}\NormalTok{, }\AttributeTok{quietly =} \ConstantTok{TRUE}\NormalTok{)) }
  \FunctionTok{install.packages}\NormalTok{(}\StringTok{"devtools"}\NormalTok{)}
\NormalTok{devtools}\SpecialCharTok{::}\FunctionTok{install\_github}\NormalTok{(}\StringTok{"aphrcwaro/cd2030.rmncah@v1.0.0"}\NormalTok{)}
\end{Highlighting}
\end{Shaded}

\textbf{Development Version}

\begin{Shaded}
\begin{Highlighting}[]

\ControlFlowTok{if}\NormalTok{ (}\SpecialCharTok{!}\FunctionTok{requireNamespace}\NormalTok{(}\StringTok{"devtools"}\NormalTok{, }\AttributeTok{quietly =} \ConstantTok{TRUE}\NormalTok{)) }
  \FunctionTok{install.packages}\NormalTok{(}\StringTok{"devtools"}\NormalTok{)}
\NormalTok{  devtools}\SpecialCharTok{::}\FunctionTok{install\_github}\NormalTok{(}\StringTok{"aphrcwaro/cd2030.rmncah"}\NormalTok{)}
\end{Highlighting}
\end{Shaded}

\subsection{Launching the Application}\label{launching-the-application}

\begin{Shaded}
\begin{Highlighting}[]

\FunctionTok{library}\NormalTok{(cd2030.rmncah)}
\FunctionTok{dashboard}\NormalTok{()}
\end{Highlighting}
\end{Shaded}

The Shiny dashboard will launch automatically in your default web
browser.

\begin{center}
\includegraphics[width=1\linewidth,height=\textheight,keepaspectratio]{pages/../images/dashboard.jpg}
\end{center}

\subsection{Alternative Installation via GitHub Desktop or Git or direct
download}\label{alternative-installation-via-github-desktop-or-git-or-direct-download}

Advanced users with GitHub accounts may prefer to clone the repository
directly. This method allows for:

\begin{itemize}
\tightlist
\item
  Version control
\item
  Contribution to the codebase
\item
  Inspection of the package structure
\item
  Once cloned, open the .Rproj file in RStudio to set the working
  directory.
\end{itemize}

\begin{center}
\includegraphics[width=1\linewidth,height=\textheight,keepaspectratio]{pages/../images/github-repo.png}
\end{center}

\begin{center}
\includegraphics[width=1\linewidth,height=\textheight,keepaspectratio]{pages/../images/project-file.png}
\end{center}

To install and run:

\begin{Shaded}
\begin{Highlighting}[]
\FunctionTok{library}\NormalTok{(cd2030.rmncah)}
\FunctionTok{dashboard}\NormalTok{()}
\end{Highlighting}
\end{Shaded}

\subsection{\texorpdfstring{\texttt{RTools} for
Windows}{RTools for Windows}}\label{rtools-for-windows}

To compile packages from source, especially development versions, RTools
must be installed. Ensure the Rtools selected corresponds and is
compatibile with your installed R version.:

\begin{itemize}
\tightlist
\item
  \textbf{R 4.2.0} → Rtools42
\item
  \textbf{R 4.4.0} → Rtools44
\item
  \textbf{R 4.5.0} → Rtools45
\end{itemize}

Rtools can be downloaded from
\href{https://cran.r-project.org/bin/windows/Rtools/}{CRAN}

\section{Data requirements}\label{sec-data-requirements}

\subsection{Datasets required}\label{datasets-required}

To run the analysis efficiently, each country team will require to have
a folder containing the following datasets:

\begin{enumerate}
\def\labelenumi{\arabic{enumi}.}
\tightlist
\item
  Health facility data (.xlsx file)
\item
  UN Estimates
\item
  Survey data
\item
  UN mortality data
\item
  FPET data
\end{enumerate}

These datasets will be provided to the country teams by the
Countdown2030 team prior to the workshop.

\subsection{Country specific analysis
parameters}\label{country-specific-analysis-parameters}

The following (\hyperref[sec-denom-analysis-setup]{parameters}) will be
required to run the analysis

\chapter{App Features}\label{app-features}

\section{Overview of App Features}\label{overview-of-app-features}

This section details the main components of the \textbf{cd2030.rmncah}
Shiny App interface: the \textbf{Title Bar}, \textbf{Sidebar}, and
\textbf{Main Panel (Body Content)}. Understanding these elements will
help you navigate and use the app more effectively.

\section{Title Bar}\label{title-bar}

The Title Bar, located at the top of the app, provides key information
controls. The elements of the Title Bar, such as ``Set Cache
Directory/Save Cache'' and ``Download Report,'' are only displayed after
data has been loaded.

\subsection{Key Elements of the Title
Bar}\label{key-elements-of-the-title-bar}

\begin{enumerate}
\def\labelenumi{\arabic{enumi}.}
\tightlist
\item
  \textbf{Contextual Information}: Displays the current analysis
  context, such as the country of analysis (e.g., ``Madagascar -
  Countdown Analysis'').
\end{enumerate}

\begin{center}
\includegraphics[width=1\linewidth,height=\textheight,keepaspectratio]{pages/../images/0-sidebar_navigation.png}
\end{center}

\begin{enumerate}
\def\labelenumi{\arabic{enumi}.}
\setcounter{enumi}{1}
\item
  \textbf{Cache Management}

  \begin{center}
  \includegraphics[width=1\linewidth,height=\textheight,keepaspectratio]{pages/../images/0-cache-setting.png}
  \end{center}

  \begin{itemize}
  \item
    \textbf{Set Cache Directory:} This button allows you to specify a
    directory for saving intermediate results and progress.
  \item
    \textbf{Save Cache:} This button allows you to save the progress of
    your analysis. This is useful for large analyses or for resuming a
    session later.
  \item
    NB:

    \begin{itemize}
    \tightlist
    \item
      The Cache button is enabled only once valid data has been uploaded
      into the App
    \item
      Changing the file path of the Cache mid analysis would require you
      to set the cache again, losing any saved progress after the last
      save.
    \end{itemize}
  \end{itemize}
\end{enumerate}

\section{Sidebar}\label{sidebar}

The Sidebar is located on the left side of the app and contains the main
navigation elements.It provides navigation through the analysis
workflow. The sections are arranged sequentially to guide you through
the process.

\subsection{Sidebar Navigation}\label{sidebar-navigation}

\begin{center}
\includegraphics[width=1\linewidth,height=\textheight,keepaspectratio]{pages/../images/0-sidebar_navigation.png}
\end{center}

\begin{enumerate}
\def\labelenumi{\arabic{enumi}.}
\item
  \textbf{Introduction}: Provides a brief overview of the app's purpose
  and functionality.
\item
  \textbf{Load Data}: Allows you to upload and manage your datasets
  (e.g.~Health Facility Datasets, Master datasets, and cache file).
\item
  \textbf{Data Quality}: Provides tools for assessing the quality of
  your data, including checks for missing values and inconsistencies.
  Subsections include:

  \begin{itemize}
  \tightlist
  \item
    Reporting Rate
  \item
    Outlier Detection
  \item
    Consistency Check
  \item
    Data Completeness
  \item
    Calculate Ratios
  \item
    Overall Score
  \end{itemize}
\end{enumerate}

\begin{center}
\includegraphics[width=1\linewidth,height=\textheight,keepaspectratio]{pages/../images/1-dqa_reporting_rates.png}
\end{center}

\begin{enumerate}
\def\labelenumi{\arabic{enumi}.}
\setcounter{enumi}{3}
\tightlist
\item
  \textbf{Analysis Set-up}: Allows you to configure the analysis by
  selecting the data sources (UN Estimates, WUENIC estimates, and survey
  data). It also allows you to map the sub-national survey and mapping
  data to the uploaded data.
\end{enumerate}

\begin{center}
\includegraphics[width=1\linewidth,height=\textheight,keepaspectratio]{pages/../images/1-dqa_analysis_setup.png}
\end{center}

\begin{enumerate}
\def\labelenumi{\arabic{enumi}.}
\setcounter{enumi}{4}
\tightlist
\item
  \textbf{National Analysis}: Provides tools for analyzing
  national-level data, including coverage and equity analysis.
\item
  \textbf{Sub-National Analysis:} For all sub-national analyses
\item
  \textbf{Mortality Analysis:} For all mortality analyses
\item
  \textbf{Service Utilization:} ForCurative health services utilization
  for sick children analyses
\item
  \textbf{Health System Performance:} Health systems progress and
  performance analysis
\end{enumerate}

\section{Main Panel(Body Content)}\label{main-panelbody-content}

The main panel is the central workspace where the analysis result,
visualizations, and interactive elements are displayed.

\subsection{Key Elements of the Main
Panel}\label{key-elements-of-the-main-panel}

\begin{enumerate}
\def\labelenumi{\arabic{enumi}.}
\tightlist
\item
  \textbf{Page Title}: Displays the current module or analysis step
  (e.g., ``Reporting Rate'').
\end{enumerate}

\begin{center}
\includegraphics[width=1\linewidth,height=\textheight,keepaspectratio]{pages/../images/1-dqa_reporting_rates.png}
\end{center}

\begin{enumerate}
\def\labelenumi{\arabic{enumi}.}
\setcounter{enumi}{1}
\tightlist
\item
  \textbf{Action Buttons}: Provides context-specific actions:

  \begin{itemize}
  \item
    \textbf{Download Report}: Helps you to download two reports:
    \textbf{Synthesis report} and the \textbf{Subnational one paper
    report}
  \item
    \textbf{Generate Report}: Downloads a report specific to the
    analysis section.

    NB: This button is green in colour and only appears in the sections
    where the user can download the section reports.
  \item
    \textbf{Get Help}: Provides context-specific help and documentation
    for the current page.
  \item
    \textbf{Download Plot / Download Data:} Allows the user to download
    the displayed output and the data associated with the output,
    respectively.
  \end{itemize}
\end{enumerate}

\chapter{Loading Health facility
data}\label{loading-health-facility-data}

\begin{center}\rule{0.5\linewidth}{0.5pt}\end{center}

This section explains how to structure and upload data into the app,
ensuring compatibility with the \textbf{Countdown Health facility data
extraction template}

Please note that should you update your data during the workshop, ensure
it is in the correct format before uploading to avoid errors.

\section{\texorpdfstring{\textbf{Supported File
Formats}}{Supported File Formats}}\label{supported-file-formats}

The app supports uploading the following file types:

\begin{itemize}
\item
  \texttt{.xls}, \texttt{.xlsx} (Excel files)- The raw health facility
  dataset in the \textbf{Countdown data template}
\item
  \texttt{.dta} (Stata files) - Master dataset downloaded from the app
  after validation/adjustment
\item
  \texttt{.rds} (R Cached datafile) - The file containing the preloaded
  dataset, user adjustments and analysis parameters that has been
  \textgreater{} saved in the Cache directory. This will be the last
  saved file.
\end{itemize}

\section{\texorpdfstring{\textbf{How to Upload
Data}}{How to Upload Data}}\label{how-to-upload-data}

\textbf{Step 1: Prepare your data file}

Ensure your data is cleaned and structured according to the
\textbf{Countdown Health facility data (HFD)} template by:

\begin{itemize}
\item
  Using the provided \textbf{HFD Standardized Template} to format your
  data correctly.
\item
  Saving the file in a supported format: \texttt{.xls}, \texttt{.xlsx},
  .\texttt{dta}, or \texttt{.rds}.
\end{itemize}

Step 2: Upload the File

\begin{enumerate}
\def\labelenumi{\arabic{enumi}.}
\item
  Navigate to the \texttt{Upload\ Data} section of the app.
\item
  \textbf{Drag and drop} your file into the upload box, or click
  \textbf{\texttt{Browse}} \textgreater{} to select it manually in your
  directory.
\item
  For subsequent re-uploads (after the initial uploading of the
  \texttt{.xls}, \texttt{.xlsx}, \texttt{.dta} files and saving your
  progress in .rds file using the \texttt{Save\ Cache} button), do not
  re-upload the \texttt{.xls}, .\texttt{xlsx}, or \texttt{.dta} but the
  saved \texttt{.rds} file if you want to retain any changes made in
  your analytical files
\item
  The app will validate your file against the \textbf{Countdown
  template}.

  \begin{itemize}
  \item
    If successful, a confirmation message will appear:
    \textbf{\emph{``Upload successful: Your file is ready for
    analysis.''}}
  \item
    If errors are detected, an error message will indicate the issue.
  \end{itemize}
\end{enumerate}

\section{\texorpdfstring{\textbf{Common Errors and How to Fix
Them}}{Common Errors and How to Fix Them}}\label{common-errors-and-how-to-fix-them}

\begin{longtable}[]{@{}
  >{\raggedright\arraybackslash}p{(\linewidth - 4\tabcolsep) * \real{0.3333}}
  >{\raggedright\arraybackslash}p{(\linewidth - 4\tabcolsep) * \real{0.3333}}
  >{\raggedright\arraybackslash}p{(\linewidth - 4\tabcolsep) * \real{0.3333}}@{}}
\toprule\noalign{}
\begin{minipage}[b]{\linewidth}\raggedright
\textbf{Error Message}
\end{minipage} & \begin{minipage}[b]{\linewidth}\raggedright
\textbf{Cause}
\end{minipage} & \begin{minipage}[b]{\linewidth}\raggedright
\textbf{Solution}
\end{minipage} \\
\midrule\noalign{}
\endhead
\bottomrule\noalign{}
\endlastfoot
``Unsupported file format'' & File type not supported & Save your file
as .xls, .xlsx, .dta, or .rds. \\
``The following required columns are missing from the data: opv1'' &
Missing essential columns in the data & Add the missing column(s) to
your dataset and ensure their values are valid. \\
``The following sheets are missing: Service\_data\_1, Service\_data\_2,
Service\_data\_3, Reporting\_completeness, Population\_data,
Admin\_data'' & Missing one or more required sheets in the file & Add
the missing sheets to your file and ensure they conform to the
template. \\
``Sheet Service\_data\_3 is empty'' & The sheet exists but contains no
data & Populate the sheet with valid data or remove the empty sheet. \\
``Key Columns''month'' missing in Service\_data\_3'' & A key column
district, year, or month, is missing from the specified sheet & Add the
missing column(s) to the sheet and ensure the data is structured
correctly. \\
``Column name month must not be duplicated. Use .name\_repair to specify
repair.'' & Duplicate column names in the dataset & Ensure all column
names are unique. Rename or remove duplicate columns. \\
\end{longtable}

\emph{\textbf{Note}: If a key column (district, year, or month) is
missing data in a row, that row will be excluded from the resulting
dataset.}

\section{\texorpdfstring{\textbf{Tips for a Successful
Upload}}{Tips for a Successful Upload}}\label{tips-for-a-successful-upload}

\begin{itemize}
\item
  Always use the latest the \textbf{Countdown Health facility data
  template} \textgreater{} template to structure your data.
\item
  Double-check column names, formats, and content before uploading.
\item
  Save your file in a supported format and ensure it is UTF-8 encoded
\end{itemize}

\part{Data Quality Assessment}

\chapter{Numerators Assessment}\label{numerators-assessment}

\section{Rationale, Approach and
implementation}\label{sec-dqa-assess-science}

\textbf{Rationale: Scientific basis for the analysis}

Routinely reported health facility data are an important data source for
health indicators at facility and population level. The data are
reported by health facilities on events such as immunizations given, or
live births attended. As with any data, quality is an issue. Data needs
to be checked to consider completeness of reporting by health
facilities, identify extreme outliers and internal consistency. A
standard reporting method for data quality allows assessment of progress
over time.

\textbf{Approach: Description of analytical steps}

The analysis of the monthly data by district for 2019-2024 is used to
assess annual data quality using the following standard indicators:
\emph{Reporting completeness}, \emph{Outlier detection},\emph{Internal
Consistency} \emph{Data completeness}, \emph{Ratios calculation},
\emph{Overall quality score}

\textbf{Implementation: Conducting analysis in the Shiny App}

\section{DQA Metrics calculation and
interpretation}\label{dqa-metrics-calculation-and-interpretation}

\subsection{Reporting Completeness}\label{sec-dqa-reporting-rate}

\begin{longtable*}[t]{llll}
\toprule
\multicolumn{4}{c}{Completeness of Monthly Facility Reporting} \\
\cmidrule(l{3pt}r{3pt}){1-4}
Indicator & Numerator & Denominator & Interpretation\\
\midrule
1a & N of monthly facility reports received & Total N of facility reports expected (usually 12 × N of facilities with the service) & Reporting rates over 90\% are good; changes in reporting completeness over time affects trend analysis\\
1b & N of districts with at least 90\% monthly reporting completeness in a year & Total N of districts & Can be used to identify districts with low reporting rates in multiple years\\
1c & N of districts with no missing values for any of the 4 forms in a year & Total N of districts & Additional indicator, is not used to compute the overall data quality score\\
\bottomrule
\multicolumn{4}{l}{\rule{0pt}{1em}Statistics for 1a and 1b are based on the mean of 4 reporting forms (ANC, delivery, immunization, OPD)}\\
\end{longtable*}

Within the Shiny App, you can select the indicator of interest, the
reporting rate cutoff point (performance Threshold)- the default value
has been set at 90\% reporting rate.

You can also choose the level of you analysis (Admin Level (Regions) or
District level). Below is an output for ANC reporting rate at the Admin
Level 1.

\begin{center}
\includegraphics[width=1\linewidth,height=\textheight,keepaspectratio]{pages/../images/1-dqa_reporting_rates.png}
\end{center}

The app will also produce a report with the results of the data quality
assessment, including the numerators and denominators for each
indicator, as well as the overall data quality score.

\subsection{Data Completeness}\label{sec-dqa-data-completeness}

\begin{center}
\includegraphics[width=1\linewidth,height=\textheight,keepaspectratio]{pages/../images/1-dqa_data_completeness.png}
\end{center}

\subsection{Outlier Detection}\label{sec-dqa-outlier}

\begin{longtable*}[t]{llll}
\toprule
\multicolumn{4}{c}{Extreme Outliers} \\
\cmidrule(l{3pt}r{3pt}){1-4}
Indicator & Numerator & Denominator & Interpretation\\
\midrule
2a & N of monthly values that are not extreme outliers in a specific year & Total N of monthly values (usually 12 * N years to be analyzed) & At least 99\% of monthly data expected not to be an extreme outlier; consider reasons\\
2b & N of districts with no extreme outliers in a specific year & Total N of districts & At least 90\% of districts should have no extreme outliers at all; consider reasons\\
\bottomrule
\multicolumn{4}{l}{\rule{0pt}{1em}Outliers are identified statistically; definitions may vary depending on methods used.}\\
\end{longtable*}

The Shiny App allows users to visually investigate which sub-national
units have outliers and for which months. This way, they can factor in
their country context in determining whether indeed it is a data quality
issue or not.

\begin{center}
\includegraphics[width=1\linewidth,height=\textheight,keepaspectratio]{pages/../images/1-dqa_outliers.png}
\end{center}

\subsection{Internal Consistency}\label{sec-dqa-consistency}

\begin{longtable*}[t]{llll}
\toprule
\multicolumn{4}{c}{Consistency of Annual Reporting} \\
\cmidrule(l{3pt}r{3pt}){1-4}
Indicator & Numerator & Denominator & Interpretation\\
\midrule
3a & N of ANC1 reported & N of penta1 reported & National ratio within an expected range (1.05 to 1.10 if survey coverage ANC1 and penta1 are the same – see below)\\
3b & N of penta1 reported & N of penta3 reported & National ratio within an expected range, based on the survey results (see below)\\
3c & N of districts with ratios within the expected range & Total N of districts & For districts there is more variation in the ratio: a wider range is considered\\
3d & N of districts within the expected range & Total N of districts & For districts there is more variation in the ratio: a wider range is considered\\
\bottomrule
\end{longtable*}

There is often an inconsistency between antenatal and immunization data,
even though we can argue that the two should be consistent. To examine
the association between ANC1 and penta1 is particularly informative.

\begin{center}
\includegraphics[width=1\linewidth,height=\textheight,keepaspectratio]{pages/../images/1-dqa_internal_consistency.png}
\end{center}

To compute and interpret indicators 3a and 3b the following
considerations need to be made:

\textbf{\emph{ANC1 to penta1 ratio}}

We can compute an expected ratio ANC1 to penta1 based on assumptions
about mortality between early to mid pregnancy and early infancy and
survey data on coverage of ANC1 and penta1 in the population:

\begin{itemize}
\item
  Consider the mortality between the first ANC visit and the first
  pentavalent vaccination.

  Assuming that ANC1 takes place at about 20 weeks or 4-5 months of
  pregnancy and penta1 at 6-8 weeks postpartum, we assume a pregnancy
  loss (abortion) after the ANC1 visit of 3\%, a stillbirth rate of 2\%,
  a twinning rate of 1.5\% and neonatal mortality rate before the penta1
  of 3\% then the difference between the numbers of ANC1 and penta1
  should be: \[1 - 0.03 -0.02 + 0.015 - 0.03 = 0.935 \] This corresponds
  with a ANC1 to penta1 ratio of \[1/0.935 = 1.07\]
\item
  Actual population coverage of ANC1 and penta1 will also need to be
  considered, using the surveys.

  The expected ratio (the number of ANC1/ number of penta1 in
  facilities) is 1.07 * (ANC1 coverage in the survey/penta1 coverage in
  the survey). If coverage for ANC1 and penta1 are the same, then the
  ratio is \[ 1.07 (1.07 \times 1/1) \]

  But if, for example, the last survey shows that ANC1 coverage was 90\%
  and penta1 coverage was 95\%, then the expected ratio becomes
  \textbackslash1.07 * (.90/.95) = 1.01*.
\item
  For the national ANC1 to penta1 ratio a range of plus or minus 0.05
  outside this computed ratio is considered acceptable. If the ratio is
  outside this range, this should be flagged, and possible explanations
  discussed.
\end{itemize}

\textbf{\emph{Penta1 to penta3 ratio}}

We can compute an expected penta1 to penta3 ratio based on the most
recent survey:

\begin{itemize}
\item
  The main factor determining the penta1 to penta3 ratio, which are
  recommended at 6 and at 14 weeks of age, is the actual drop-out rate
  between penta1 and penta3, as mortality plays a limited role.
\item
  Population coverage rates from the latest survey are used to determine
  the expected penta1 to penta3 ratio in the facility data. For
  instance, if penta1 coverage is 95\% and penta3 coverage is 85\%, we
  expect that ratio to be 0.95/0.85 = 1.12.
\item
  Also, here a range of plus or minus 0.05 is considered acceptable for
  the assessment of the facility data.
\end{itemize}

The Figure below produced in the ShinyApp shows the ratios for all six
years (2019-2024) and the expected value. It will be important to
reflect on large differences (e.g., more than 0.10 or 10\%.).

\begin{center}
\includegraphics[width=1\linewidth,height=\textheight,keepaspectratio]{pages/../images/1-dqa_ratios.png}
\end{center}

\subsection{Overall DQA Score}\label{sec-dqa-overall-score}

In this section, the overall DQA score is calculated by considering the
mean scores within some preselected data quality metrics as below.

\begin{longtable*}[t]{llll}
\toprule
\multicolumn{4}{c}{Summary of Performance} \\
\cmidrule(l{3pt}r{3pt}){1-4}
Indicator & Numerator & Denominator & Interpretation\\
\midrule
 &  &  & Annual data quality score (mean 1a, 1b, 2a, 2b, 3c, 3d)\\
\bottomrule
\end{longtable*}

The app will automatically compute the overall data quality score based
on the metrics described above. It is in this section that the first
section report is downloaded capturing the key DQA indicators as shown
below.

\begin{center}
\includegraphics[width=1\linewidth,height=\textheight,keepaspectratio]{pages/../images/1-dqa_overall_score.png}
\end{center}

\subsection{\texorpdfstring{\textbf{DQA Report and
interpretation}}{DQA Report and interpretation}}\label{dqa-report-and-interpretation}

The first report - DQA - is generated using the
\texttt{Generate\ Report} button.

The interpretation for this section should seek to answer the following
questions:

\begin{itemize}
\tightlist
\item
  Is there a data quality pattern by year for which there is an
  explanation? (include the explanation)
\item
  Are there certain regions or other sub-national units that are
  particularly problematic?
\item
  Are there certain reporting forms or services (e.g., antenatal care,
  labour and delivery, immunization) that are problematic?
\item
  Is there good consistency between reported numbers of ANC1 and penta1?
\end{itemize}

¹ An extreme outlier is defined as a monthly value that is 5 times the
median absolute deviation (MAD) from the monthly median value for a
particular year.

\chapter{Numerator adjustments}\label{numerator-adjustments}

\section{Rationale, Approach and
implementation}\label{sec-dqa-adjust-science}

\textbf{Rationale: Scientific basis for the analysis}

The completeness of reporting may affect the analysis, especially if
completeness is low or varies between years. Extreme outliers, such as
an accidental extra zero in a number, can have a large impact,
especially on subnational numbers. Following the assessments, several
steps are necessary to obtain a clean data set for analysis. This
implies adjusting for incomplete reporting and correcting for extreme
outliers.

\textbf{Approach: Description of analytical steps}

If we do not consider reporting completeness that means we assume all
non-reporting facilities provided zero services, which is not likely to
be true. Adjustments depend on how much services (e.g., pregnancy care,
vaccinations) were provided at non-reporting facilities compared to
those that reported. The adjustment factor k - defined as the ratio of
the volume of services provided by non-reporting facilities to the
volume of services provided by reporting facilities - is used to adjust
the reported numbers for incomplete reporting.

\section{Numerator Adjustment}\label{sec-numerator-adjustment}

To account for incomplete reporting, the reported number of events can
be adjusted using completeness and facility reporting ratio, with the
following formula:

\[
N_{\text{adjusted}} = N_{\text{reported}} + N_{\text{reported}} \times \left( \frac{1}{c} - 1 \right) \times k
\]

\textbf{Where:}

\begin{itemize}
\tightlist
\item
  \(N_{\text{adjusted}}\): Total number of events adjusted for
  incomplete reporting
\item
  \(N_{\text{reported}}\): Number of events reported
\item
  \(( c )\): Reporting completeness (e.g., proportion of facilities that
  reported)
\item
  \(( k )\): Adjustment factor to account for lower service volume in
  non-reporting facilities
\end{itemize}

As a default value, we use \(k=0.25\) , which means the non-reporting
facilities provided services but only at a volume which was a quarter of
the reporting health facilities.

The factor k can be different for different services. For instance, if
private facility reporting is poor but they are in the national system
and they provide a considerable number of deliveries, k maybe greater
than 0.25 or even as high as 1.0.

The following k-values are used depending on the reporting used to
adjust the reported numbers for incomplete reporting:

\begin{itemize}
\tightlist
\item
  \(k=0\) - No services in non-reporting facilities (default k-value)
\item
  \(k=0.25\) \textbf{\emph{-}} Some services, but much lower than
  reporting facilities
\item
  \(k=0.50\) \textbf{\emph{-}} Half the rate compared to reporting
  facilities
\item
  \(k=0.75\) \textbf{\emph{-}} Nearly as much as reporting facilities
\item
  \(k= 1.0\) - Same rate of services as reporting facilities
\end{itemize}

If the facility reporting rate is below 75\%, it becomes more difficult
to impute district data. Therefore, no adjustments are made if reporting
is lower than 75\%. In that case, further analysis to determine coverage
with the facility data is not considered sufficiently reliable.

Extreme outliers, as defined in the previous section, will be corrected
by imputing the median monthly value of the same year. The table
summarizes the adjustments.

\textbf{Table 2: Summary of adjustments made to the raw health facility
data in preparation of a clean data set for the endline analysis}

\begin{longtable}[]{@{}
  >{\raggedright\arraybackslash}p{(\linewidth - 4\tabcolsep) * \real{0.3333}}
  >{\raggedright\arraybackslash}p{(\linewidth - 4\tabcolsep) * \real{0.3333}}
  >{\raggedright\arraybackslash}p{(\linewidth - 4\tabcolsep) * \real{0.3333}}@{}}
\toprule\noalign{}
\begin{minipage}[b]{\linewidth}\raggedright
\textbf{Problem}
\end{minipage} & \begin{minipage}[b]{\linewidth}\raggedright
\textbf{Action}
\end{minipage} & \begin{minipage}[b]{\linewidth}\raggedright
\textbf{Adjustment}
\end{minipage} \\
\midrule\noalign{}
\endhead
\bottomrule\noalign{}
\endlastfoot
Low reporting rates: identifying low rates that were adjusted & If below
75\% (default), data were imputed & Median monthly value for the
district year was imputed for the month with low reporting \\
Incomplete reporting by districts, variable over time, affecting trend
assessment & If reporting rates were \textgreater=75\% and \textless=
100\% default), an assumption was made about the volume of services
provided by the non-reporting facilities & Adjustment factor k value was
used to adjust for incomplete reporting k default value 0.25 (replace if
different value used; state if used for all reporting forms or different
k factors between forms) \\
Extreme outliers can greatly affect coverage trend assessments & If a
monthly value was greater or smaller than 5 times the median absolute
deviation (MAD) from district monthly median value, an adjustment was
made & Extreme monthly outliers are corrected and given the district
median value for the same year \\
Missing values & If there is a missing value, data were imputed &
District median monthly value for the year was imputed for the month
with missing value \\
\end{longtable}

\section{Implementation: Conducting analysis in the Shiny
App}\label{sec-dqa-adjust-implement}

\begin{itemize}
\item
  The outputs for this analysis can be obtained through the \textbf{Data
  adjustment} section of the Shiny App.
\item
  Select the \(k-factor\) that is deemed appropriate given your country
  context.

  \begin{center}
  \includegraphics[width=1\linewidth,height=\textheight,keepaspectratio]{pages/../images/1-dqa-k_factor.png}
  \end{center}
\end{itemize}

\subsection{Remove Years}\label{remove-years}

The \textbf{\emph{\texttt{Remove\ Years}}} section allows you to remove
any years with no data or whose data has been determined to not be fit
for analysis by the country team due to quality issues.

It is recommended to be used in instances where the data quality is
poor, or the country teams have sufficient, contextual information that
renders the said year(s)' data unreliable/inaccurate.This will ensure
that the analysis is based on reliable data.

\begin{center}
\includegraphics[width=1\linewidth,height=\textheight,keepaspectratio]{pages/../images/1-dqa_remove_years.png}
\end{center}

\subsection{Adjustment outputs}\label{sec-dqa-adjust-outputs}

\begin{itemize}
\item
  This Section produces a master adjusted dataset \texttt{(.dta)} ready
  for analysis and plots for different indicators showing data
  adjustment changes. In the image below, we can see the adjustment
  output after making no adjustment to the data \((k=0)\).

  \begin{center}
  \includegraphics[width=1\linewidth,height=\textheight,keepaspectratio]{pages/../images/1-dqa_adjustments.png}
  \end{center}
\end{itemize}

\subsection{\texorpdfstring{\textbf{Report}}{Report}}\label{sec-dqa-adjust-report}

\begin{itemize}
\tightlist
\item
  The interpretation should include the selected adjustment factor
  (factor k) that was used to adjust for incomplete reporting (if
  necessary, by service). If the default factor is used, then report
  this and explain what this means for the reader.
\item
  Report the percent change that the adjustment made in reported numbers
  of institutional deliveries and in penta1 (average of the 6-year
  period);
\item
  You may want to :

  \begin{itemize}
  \tightlist
  \item
    Highlight the year with the greatest impact of the adjustment if
    there is one;
  \item
    Interpret if the impact of the adjustment on coverage rates is large
    or small;
  \item
    Make the same description and interpretations for penta1
    vaccinations.
  \end{itemize}
\end{itemize}

\chapter{Denominator Assessment and
selection}\label{sec-denom-assessment}

\section{Rationale, Approach and
Implementation}\label{rationale-approach-and-implementation}

\textbf{Rationale: Scientific basis for the analysis}

Service coverage is defined as the population who received the service
divided by the population who need the services (also referred to as
target population). The numerators of coverage statistics (e.g., number
of live births in health facilities) are derived from the health
facility data and need to be adjusted as shown in the previous section.
The denominator of the coverage statistics (e.g., number of live births
in the population) needs to be estimated for national and sub-national
levels (regions and districts).

\textbf{Approach: Description of analytical steps}

The objective of the health facility denominator analysis is twofold:

\begin{enumerate}
\def\labelenumi{\arabic{enumi}.}
\tightlist
\item
  First, we assess the quality of the population projections in DHIS2 by
  comparison with the UN projections and the internal consistency.
\item
  Then, we assess the performance of multiple denominators options for
  the computation of population-based service coverage indicators from
  the health facility data.
\end{enumerate}

This should lead to a final decision on denominators that are used for
the analyses of population-based coverage indicators based on health
facility data.

Each indicator has its own denominator, as shown in the table below.

\textbf{Table 3: Selected indicators with numerators and denominators}

\begin{longtable}[]{@{}
  >{\raggedright\arraybackslash}p{(\linewidth - 4\tabcolsep) * \real{0.3333}}
  >{\raggedright\arraybackslash}p{(\linewidth - 4\tabcolsep) * \real{0.3333}}
  >{\raggedright\arraybackslash}p{(\linewidth - 4\tabcolsep) * \real{0.3333}}@{}}
\toprule\noalign{}
\begin{minipage}[b]{\linewidth}\raggedright
\textbf{Indicators}
\end{minipage} & \begin{minipage}[b]{\linewidth}\raggedright
\textbf{Numerator}
\end{minipage} & \begin{minipage}[b]{\linewidth}\raggedright
\textbf{Denominator}
\end{minipage} \\
\midrule\noalign{}
\endhead
\bottomrule\noalign{}
\endlastfoot
\textbf{SERVICE UTILIZATION} & & \\
Outpatient visits, children under 5, per year (N) & N of OPD visits for
under-5 & Total mid-year population under 5 \\
Inpatient admissions, children under 5, per year (N) & N of admissions
for under-5 & Total mid-year population under 5 \\
\textbf{PREVENTIVE INTERVENTIONS} & & \\
\% of pregnant women with 4 antenatal care visits & N of women with ANC
4th visit & Total N of pregnant women in the whole population \\
\% of live births in health facilities & N of live births in health
facilities & Total N of live births in the whole population \\
\% of infants receiving 3 doses of pentavalent vaccine & N of infants
receiving 3 doses & N of infants eligible for 3 doses of the vaccine \\
\textbf{CURATIVE INTERVENTIONS} & & \\
\% of children under 5 with malaria who receive ACT & N of children
under 5 with malaria receiving ACT & Total N of children who had malaria
in the last year \\
\% of deliveries that were by C-section (population) & N of C-sections
reported & Total N of deliveries in the population \\
\% of deliveries that were by C-section (institutional) & N of
C-sections reported & Total N of deliveries in health facilities \\
\textbf{MORTALITY} & & \\
Institutional maternal mortality ratio & N of maternal deaths in health
facilities & Total number of live births in health facilities \\
Stillbirth rate & N of stillbirths in health facilities & Total N of
births in health facilities \\
Neonatal mortality before discharge & N of neonatal deaths before
discharge (after birth) & Total N of live births in the health
facilities \\
\textbf{FAMILY PLANNING (FP)} & & \\
Ratio FP visits to women of reproductive age & N of FP new and revisits
& Total N of women 15-49 years \\
Estimated modern use of contraceptives & Couple years of protection &
Total N of women 15-49 years \\
FP coverage (demand satisfied) & N of women using modern methods & Total
N of women in need of FP \\
\end{longtable}

\subsection{Part 1: Quality assessment of population projections in
DHIS2}\label{part-1-quality-assessment-of-population-projections-in-dhis2}

In the first part, we assess the \textbf{quality of the DHIS2 population
projections} at national level:

\begin{itemize}
\tightlist
\item
  \textbf{Check the internal consistency of the DHIS2 population growth
  over time:}
\end{itemize}

Compute the population growth rate:

\[
\text{Pop growth} = \frac{\text{Pop}_{2024} - \text{Pop}_{2023}}{\text{Pop}_{2023}}, 
\]then the crude birth rates (defined as the number of live births in
DHIS2 projections per 1,000 population). We expect both rates to be
consistent over time (e.g., less than 2 per 1,000 difference between
years).

\begin{itemize}
\tightlist
\item
  \textbf{Compare the population data in DHIS2 with the UN population
  projections at the national level:}
\end{itemize}

Differences may occur, but large discrepancies suggest issues with DHIS2
population projections. The comparison is done for four indicators.
Abnormal values are flagged:

\begin{enumerate}
\def\labelenumi{\arabic{enumi}.}
\item
  \textbf{Population size:}\\
  A relative difference between DHIS2 and UN-projected population size
  greater than 5\% should be considered a data quality issue with the
  DHIS2 projections.

  \$ \text{Relative difference} = \left\textbar{}
  \frac{\text{Pop}_{\text{DHIS2}} - \text{pop}_{\text{UN}}}{\text{Pop}_{\text{UN}}}
  \right\textbar{} \$
\item
  \textbf{Population growth during 2023--2024:}\\
  Annual growth is computed using the natural logarithm as below.
\end{enumerate}

\textbf{Interpretation:} A difference greater than 0.3\% (absolute)
between DHIS2 and UN estimates is concerning.

\$ \ln \left( \frac{\text{Pop}_{2024}}{\text{Pop}_{2023}} \right) \$

\begin{enumerate}
\def\labelenumi{\arabic{enumi}.}
\setcounter{enumi}{2}
\tightlist
\item
  \textbf{Crude birth rate (CBR):}
\end{enumerate}

Defined as the number of live births per 1,000 population. We compare
the DHIS2 CBR for 2023 with the UN estimate of the CBR for the same
year.

\textbf{Interpretation:} If the difference is greater than 5 per 1,000
population, the DHIS2 may be problematic.

DHIS2 CBR is defined as:

\$ CBR =
\frac{\text{Total Projected livebirths}}{\text{Total Population}}
\times 1000 \$

\textbf{Interpretation:} A difference greater than 5 per 1,000
population compared to the UN estimate suggests a data quality issue.

\begin{enumerate}
\def\labelenumi{\arabic{enumi}.}
\setcounter{enumi}{3}
\item
  \textbf{Crude death rate (CDR):}

  Defined as the number of deaths per 1,000 population. It is the
  difference between the CBR and the population growth if thereis no
  major out- or in-migration.

  We compute the CDR from the DHIS2 projections as shown in the equation
  below

  \textbf{Interpretation:} A \textbf{negative CDR} or a \textbf{CDR
  \textless{} 5 per 1,000} indicates inconsistency and potential data
  problems.
\end{enumerate}

\(CDR = CBR - \text{Population growth rate} ; \text{(both expressed per 1,000 population )}\)

\subsection{Part 2: Denominator
Selection}\label{part-2-denominator-selection}

The second part is to select the best performing denominator for
population coverage estimates with facility data. First, we compare the
results for different denominators at the national level.

For the national level, we evaluate 5 denominator methods. The first two
are projection methods (DHIS2 estimate and UN estimate), and the
additional three are facility data- based methods (
\textbf{\emph{ANC1-derived}}, \textbf{\emph{penta1-derived}} and
\textbf{penta1-derived adjusted for population growth} denominators).

\textbf{Note:}

\begin{itemize}
\tightlist
\item
  For the sub-national level, no UN projections are available, so we
  will use 4 methods only.
\end{itemize}

The maternal and newborn health denominators are closely related and can
be computed from each other by making assumptions.

Starting with pregnancies, the number of live births is closely
associated with the number of pregnancies, which are usually identified
within the health system at the first antenatal visit which in most
countries is around 4-5 months of pregnancy (according to the surveys).
Country specific values are preferred where available and can be
obtained from the
\href{https://data.unicef.org/topic/child-survival/neonatal-mortality/}{WHO
website}. The default global assumptions are as follows:

\begin{itemize}
\tightlist
\item
  Pregnancy loss between 4 and 7 months (28 weeks of pregnancy): 3\%.
\item
  Stillbirths or pregnancy loss between 28 weeks and birth: 2\%.
\item
  Twinning rate: 1.5\%.
\end{itemize}

These first three steps give the number of live births computed from
pregnancies.

\begin{itemize}
\tightlist
\item
  Neonatal mortality: 3\% (or 30 per 1,000 live births).
\item
  Post neonatal mortality (between 1-11 months): 2.4\% (or 24 per 1,000
  live births).
\end{itemize}

The selection of the best performing denominator method is based on a
comparison of the performance of the DHIS2 projection and facility-data
derived methods for two indicators: \textbf{\emph{institutional live
births}} and \textbf{\emph{penta3}}. The gold standard is the population
coverage rates from a recent survey, for a year as close as possible.
The absolute difference between survey and facility-based coverage at
national and sub-national levels is used to select the best performing
indicator. (This can also be expressed as the number of standard errors
from the survey value but this requires including the standard errors
from the surveys - the results will be the same).

\subsection{Denominators derived from population
estimates}\label{denominators-derived-from-population-estimates}

Here we will use the estimated live births from two sources as mentioned
above;

\begin{itemize}
\item
  The UN population projections
\item
  Population projections from DHIS2 database
\end{itemize}

\subsection{Denominators derived from the routine health facility
data}\label{denominators-derived-from-the-routine-health-facility-data}

The basic idea is that if the coverage of an indicator is high (e.g.,
over 90\%), then the number of events reported by health facilities has
to be close to the target population. In other words, the denominators
or target population can be derived from the numbers in DHIS2. The best
candidate indicators for this approach are ANC1 and DPT/penta1 (BCG also
possible in some countries if re-vaccinations are recorded separately).

This approach requires the following:

A recent population-based survey is used to obtain an estimate of
population level coverage of ANC1 or penta1. For example, ANC1 coverage
is 95\% of pregnant women.

The DHIS2 data on the number of ANC1 and penta1 visits need to be
considered complete and accurate (after adjustments / cleaning the
data). For example, 100,000 ANC1 visits were reported.

If this is the case, then we only need to add the percentage that has
not used the services (according to the survey results) to get the
target population. For example, if ANC1 coverage from survey is 95\% and
the number of ANC1 visits from DHIS2 for the year is 100,000, the total
number of pregnant women is: \[ 100,000 / 0.95 = 105,263 \]

The same approach can be used for DPT1 or penta1. The survey coverage is
the percent of children 12-23 months who received DPT1/penta1, the
facility data are the number of infants who received DPT1/penta1
vaccination. For example, if survey coverage is 92\%, and there were
100,000 vaccinations given, then the
\[ \text{Denominator}= 100,000 / 0.92 = 108,696. \]

The number of live births can be obtained from ANC1 and DPT1 by making
assumptions about pregnancy loss (abortion after the first ANC visit,
stillbirths), twinning rates, and neonatal mortality. These steps are
shown in the Figure below. !{[}{]} (images/1-dqa\_denom\_adjustment.png)

\textbf{An example of ANC1:}

\begin{itemize}
\item
  Above we computed 105,263 pregnant women in the population
\item
  at 3\% abortion, this implies 105,263 * (1-.0.03) = 102,454 deliveries
\item
  at 1.5\% twinning rate this implies 102,454 / (1-(0.015/2)) = 103,229
  births
\item
  at 2\% stillbirth rate this implies 103,229 * (1-0.02) = 101,164 live
  births
\item
  at 3\% neonatal mortality this implies 101,164 * (1-0.03) = 98,129
  children eligible for DPT1/penta1.
\end{itemize}

\section{Implementation in the Shiny
App}\label{implementation-in-the-shiny-app}

\textbf{Shiny App}

To get outputs for this analysis first, one need to set up their
analysis by inputting key information at the \textbf{Analysis Set up}
section in the Shiny App.

\subsection{Analysis set up}\label{sec-denom-analysis-setup}

The parameters required are as shown in the figure below and they are:

\begin{enumerate}
\def\labelenumi{\arabic{enumi}.}
\tightlist
\item
  \textbf{National mortality rates - based on the most recent survey}
\end{enumerate}

\begin{center}
\includegraphics[width=1\linewidth,height=\textheight,keepaspectratio]{pages/../images/1-dqa_denom_rates.png}
\end{center}

\begin{enumerate}
\def\labelenumi{\arabic{enumi}.}
\item
  \textbf{ANC1 Survey(\%)} - Percentage of women aged 15--49 with a live
  birth who received at least one antenatal care (ANC) visit during
  their last pregnancy. (Example: If 940 out of 1,000 women had at least
  one ANC visit, the ANC1 coverage is 94\%.)
\item
  \textbf{Pregnancy Loss (proportion)} - Proportion of pregnancies that
  ended in miscarriage, abortion, or stillbirth out of all reported
  pregnancy outcomes. (Example: If 80 of 1,000 pregnancies ended in
  loss, the proportion is 0.08)
\item
  \textbf{Twin Births (proportion)} - Proportion of live births that are
  twins out of all live births. (Example: If 30 out of 1,500 live births
  were twins (i.e., 15 twin pairs), the proportion is 0.02)
\item
  \textbf{Neonatal Mortality (proportion)} - Proportion of live births
  that die within the first 28 days of life out of all live births.
  (Example: If 20 out of 1,000 live births died within the first 28
  days, the proportion is 0.02)
\item
  \textbf{Post Neonatal Mortality (proportion)} - Proportion of live
  births that die between 28 days and 11 months of life out of all live
  births. (Example: If 15 out of 1,000 live births died between 28 days
  and 11 months, the proportion is 0.015)
\item
  \textbf{Stillbirth Rate (proportion)} - Proportion of pregnancies that
  ended in stillbirth out of all reported pregnancy outcomes. (Example:
  If 10 out of 1,000 pregnancies ended in stillbirth, the proportion is
  0.01)
\item
  \textbf{Penta 1 Survey (\%)} - Percentage of infants aged 12--23
  months who received the first dose of pentavalent vaccine.(Example: If
  920 out of 1,000 infants received the first dose, the Penta 1 coverage
  is 92\%.)
\end{enumerate}

\begin{itemize}
\tightlist
\item
  \textbf{Survey coverage based percentages (based on the most recent
  survey)}
\end{itemize}

\begin{center}
\includegraphics[width=1\linewidth,height=\textheight,keepaspectratio]{pages/../images/1-dqa_survey_setup.png}
\end{center}

\begin{enumerate}
\def\labelenumi{\arabic{enumi}.}
\item
  \textbf{ANC4} -Percentage of women aged 15--49 with a live birth who
  received four or more antenatal care visits during their last
  pregnancy. (Example: If 680 out of 1,000 women had ≥4 ANC visits, ANC4
  coverage is 68\%.)
\item
  \textbf{Institutional delivery} - Percentage of live births that
  occurred in a health facility. (Example: If 850 out of 1,000 live
  births were in a health facility, the institutional delivery coverage
  is 85\%.)
\item
  \textbf{Low birth weight} - Percentage of live births that were low
  birth weight (less than 2,500 grams). (Example: If 100 out of 1,000
  live births were low birth weight, the low birth weight coverage is
  10\%.)
\item
  \textbf{Caesarean section} - Percentage of live births that were
  delivered by caesarean section. (Example: If 150 out of 1,000 live
  births were by C-section, the C-section coverage is 15\%.)
\item
  \textbf{Penta 3} - Percentage of infants aged 12--23 months who
  received the third dose of pentavalent vaccine. (Example: If 850 out
  of 1,000 infants received the third dose, the Penta 3 coverage is
  85\%.)
\item
  \textbf{Measles1} - Percentage of infants aged 12--23 months who
  received the first dose of measles vaccine. (Example: If 900 out of
  1,000 infants received the first dose, the Measles1 coverage is 90\%.)
\item
  \textbf{BCG} - Percentage of infants aged 12--23 months who received
  the Bacillus Calmette-Guérin (BCG) vaccine. (Example: If 950 out of
  1,000 infants received the BCG vaccine, the BCG coverage is 95\%.)
\item
  \textbf{Vaccines survey year} - The calendar year in which the survey
  was conducted and from which vaccine coverage estimates (e.g., BCG,
  Penta1/3, Measles1) are drawn. (Example: If the survey was conducted
  in 2022, that is the vaccine survey year).
\item
  \textbf{Survey data start year} - The calendar year from which the
  survey data starts. This is used to determine the time period for
  which the survey data is relevant. Ideally, it is recommended to use
  the two most recent surveys for this analysis.(Example: If a country
  has survey data for 2008, 2013, 2018 and 2023; the Survey data start
  should be 2018 (two most recent)).
\end{enumerate}

\begin{itemize}
\item
  \textbf{Survey datasets}

  In addition to setting these parameters, you will be required to
  upload the following survey datasets (in addition to the health
  facility data loaded at the beginning of the analysis session).

  \begin{center}
  \includegraphics[width=1\linewidth,height=\textheight,keepaspectratio]{pages/../images/1-dqa-survey_upload.png}
  \end{center}

  \begin{enumerate}
  \def\labelenumi{\arabic{enumi}.}
  \tightlist
  \item
    UN Estimates data
  \item
    UN Mortality data
  \item
    WUENIC estimates data
  \item
    Survey data (uploaded as a folder)
  \item
    FPET data
  \end{enumerate}
\end{itemize}

\subsection{Denominator assessment}\label{denominator-assessment}

The first part is to assess the accuracy and consistency of the
projected population numbers in the DHIS-2 by comparing them to external
sources.

\textbf{Interpretation:}

The interpretation should focus on the extent to which the DHIS2
projections are considered robust which is the case when:

\begin{itemize}
\tightlist
\item
  The DHIS2 total population projection is consistent over time with
  regular population growth
\item
  The DHIS2 total live birth projection is consistent over time (regular
  trend)
\item
  The projected numbers of total population and live births are close to
  the UN population projection
\item
  The DHIS2 population projections are consistent with UN estimates for
  crude birth rate and crude death rate.
\end{itemize}

The second part is to compare results from the different methods - both
at the national and sub-national levels.

\begin{center}
\includegraphics[width=1\linewidth,height=\textheight,keepaspectratio]{pages/../images/1-dqa_dhis2_comparison.png}
\end{center}

\textbf{Interpretation:}

The interpretation should focus on the extent to which the DHIS2
projections are considered robust which is the case when:

\begin{itemize}
\tightlist
\item
  The DHIS2 total population projection is consistent over time with
  regular population growth
\item
  The DHIS2 total live birth projection is consistent over time (regular
  trend)
\item
  The projected numbers of total population and live births are close to
  the UN population projection
\item
  The DHIS2 population projections are consistent with UN estimates for
  crude birth rate and crude death rate.
\end{itemize}

The second part is to compare results from the different methods - both
at the national and sub-national levels

\subsection{Denominator selection}\label{sec-denom-selection}

The final step is to select the best performing denominator for the
coverage analyses with health facility data. The results on the national
gap and the median sub-national gap should be taken into account to make
that choice.

\begin{center}
\includegraphics[width=1\linewidth,height=\textheight,keepaspectratio]{pages/../images/1-dqa_denom_selection.png}
\end{center}

The best methods have the smallest gaps with the survey results.

\textbf{Note}

\begin{itemize}
\tightlist
\item
  Ideally, one method is selected but it is also possible to select one
  denominator method for the MNH coverage indicators (ANC, delivery,
  PNC) and another method for the immunization coverage analyses (see
  highlighted section in the graph above). It will be important to
  clearly state the chosen denominator. Please ensure this information
  is saved within your cached \texttt{.RDS} file.
\item
  The selected denominators (for both maternal and immunization
  indicators) will be used in the subsequent analysis. The highlighted
  tabs are the only places you can change your chosen denominator(s)
\end{itemize}

\textbf{Interpretation:}

The interpretation should describe, based on the graphs:

\begin{itemize}
\tightlist
\item
  Which denominator methods performed best at the national level for the
  two indicators?
\item
  Which denominator performed best at the sub-national level for the two
  indicators?
\item
  What selection is made for the indicators in the coverage analyses?
\end{itemize}

\part{National Analyses}

\chapter{Coverage}\label{coverage}

\section{Rationale, Approach and
implementation}\label{rationale-approach-and-implementation-1}

\textbf{Rationale: Scientific basis for the analysis}

Coverage of interventions is a critical and direct output of health
systems. Regular tracking of coverage at national and sub-national
levels has become the mainstay of monitoring progress in national health
plans and international initiatives. Reproductive, maternal, newborn,
child and adolescent health indicators with targets are the most common
indicators of national health plans and global monitoring. Both health
facility data and household surveys can provide coverage statistics, and
an integrated analytical approach is desirable.

\textbf{Approach: Description of analytical steps}

Many coverage indicators can be estimated in both surveys and from
health facility data. Both are critical pieces of information and need
to be considered in conjunction with each other. The Table below lists
the indicators, for antenatal and delivery care and child immunization,
considered in the health facility data analysis in the workshop,
including the variable names (in R) in the first column.

\begin{longtable}[]{@{}
  >{\raggedright\arraybackslash}p{(\linewidth - 6\tabcolsep) * \real{0.0889}}
  >{\raggedright\arraybackslash}p{(\linewidth - 6\tabcolsep) * \real{0.4296}}
  >{\raggedright\arraybackslash}p{(\linewidth - 6\tabcolsep) * \real{0.2333}}
  >{\raggedright\arraybackslash}p{(\linewidth - 6\tabcolsep) * \real{0.2407}}@{}}
\toprule\noalign{}
\begin{minipage}[b]{\linewidth}\raggedright
\textbf{Name (in R)}
\end{minipage} & \begin{minipage}[b]{\linewidth}\raggedright
\textbf{Indicator title}
\end{minipage} & \begin{minipage}[b]{\linewidth}\raggedright
\textbf{Survey denominator}
\end{minipage} & \begin{minipage}[b]{\linewidth}\raggedright
\textbf{Facility data denominator}
\end{minipage} \\
\midrule\noalign{}
\endhead
\bottomrule\noalign{}
\endlastfoot
\textbf{Antenatal care} & & & \\
anc1 & Antenatal care at least one or more visits among all pregnant
women (\%) & Women aged 15-49 years with a live birth in the last 2
years & Estimated total pregnant women in the population \\
anc\_1 trimester & Antenatal care 1+ visits in 1st trimester of
pregnancy among all pregnant women (\%) & Women aged 15-49 years with a
live birth in the last 2 years & Estimated total pregnant women in the
population \\
anc4 & Antenatal care 4+ visits among all pregnant women (\%) & Women
aged 15-49 years with a live birth in the last 2 years & Estimated total
pregnant women in the population \\
ipt2

ipt3 & Intermittent preventive therapy for malaria - 2nd dose /
3\textsuperscript{rd} dose during pregnancy among all pregnant women
(\%) & Women aged 15-49 years with a live birth in the last 2 years &
Estimated total pregnant women in the population \\
ifa90 & Iron folic acid supplementation given (90 days supply) during
pregnancy among all pregnant women (\%) & Women aged 15-49 years with a
live birth in the last 2 years & Estimated total pregnant women in the
population \\
hiv\_test & HIV test conducted among pregnant women (\%) & Women aged
15-49 years with a live birth in the last 2 years & Estimated total
pregnant women in the population \\
syph ilis\_test & Syphilis test conducted among pregnant women (\%) &
Women aged 15-49 years with a live birth in the last 2 years & Estimated
total pregnant women in the population \\
\textbf{Delivery} \textbf{care} & & & \\
instd eliveries & Institutional live births among all live births & Live
births in the last 2 years & Estimated livebirths as denominator. \\
sba & Skilled birth attendance & Skilled birth attendance & Estimated
livebirths as denominator. \\
low\_bweight & Low birth weight (below 2500 grams) among all live births
& Low birth weight (below 2500 grams) & Estimated livebirths as
denominator. \\
c-section & C-section among all live births & C-section & Estimated
livebirths as denominator. \\
pnc48h & Postnatal care within 48 hours & Postnatal care within 48 hours
& Estimated livebirths as denominator. \\
\textbf{Immunization} & & & \\
bcg & BCG vaccination to infants & Children 12-23 months & N of
surviving infants (beyond neonatal period) \\
penta1 & Penta vaccine - 1st dose to infants & Children 12-23 months & N
of surviving infants (beyond neonatal period) \\
penta3 & Penta vaccine - 3rd dose to infants & Children 12-23 months & N
of surviving infants (beyond neonatal period) \\
measles1 & Measles vaccine - 1st dose (to infants) & Children 12-23
months & N of surviving infants (beyond post-NN period) \\
measles2 & Measles vaccine - 1st dose (older children) & Children 24-35
months & N of surviving infants (beyond postneonatal period) (or age
1) \\
\end{longtable}

The facility data can be used to generate annual coverage estimates, and
the coverage results should be compared and interpreted alongside the
results from recent surveys. The analysis results will include both
coverage estimates.

\textbf{Implementation: Conducting analysis in the Shiny App}

\subsection{Antenatal care (ANC)}\label{sec-nat-cov-anc}

Most countries have at least one ANC indicator with a target in the
national plan. The global ENAP/EPMM coverage targets for 2025 are:
globally, at least 90\% of pregnant women with 4 or more ANC care
visits, and 90\% of countries with at least 70\% coverage. There are
several ANC indicators that capture:

\begin{itemize}
\item
  \emph{contact with health services} during pregnancy (ANC
  1\textsuperscript{st} visit, ANC 4 or more visits, ANC first visit in
  first trimester). ANC1 is often considered an indicator of basic
  access to health services. It is high in most countries, and in many
  instances, the numbers of ANC1 visits in the routine health facility
  data can provide a better denominator for the ANC and delivery
  indicators than population projections (see section 2 on
  denominators).
\item
  \emph{contents of services} provided (intermittent preventive therapy
  (IPT2 or IPT3) against malaria, HIV testing, syphilis testing and
  iron-folic acid (IFA) supplementation (at least 90 tablets given to
  the pregnant woman)). Some countries will not have policies for all of
  these diagnostic or therapeutic interventions during pregnancy (e.g.,
  no IPT if no malaria risk).
\end{itemize}

For most indicators, the surveys also provide coverage estimates for the
national level, with 95\% confidence intervals. For most coverage
indicators the data refer to a period before the survey: e.g.,
institutional birth coverage for the live births in the two years
preceding the survey. This means that the midpoint of the coverage
estimate lies one year before the survey.

An example of two graphs for ANC, based on facility and survey data, is
shown below, showing both good concordance between the facility-based
and survey results for coverage of ANC first visit in the first
trimester and poor concordance in case of ANC4 visits. In the latter
case, it is evident that ANC4 is overreported in the facility data as
coverage is much higher than the survey and unlikely to be high (over
90\% during 2021-23 and even 101\% in 2023). Poor recording and
reporting of ANC 4\textsuperscript{th} visits in the DHIS2 data is
likely the main cause.

Sometimes, an indicator may reach an unlikely high coverage at the
national level, say over 125\%. This may be because the data quality of
the numerator of the coverage indicator is poor, the denominator is
wrong, or the intervention is given and recorded more than once during
pregnancy. An example is IFA supplementation. In that case, the
computation of coverage is not useful. It is better to express it
differently. For instance, if coverage is 200\%, it is better to compute
the average number of courses of 90 IFA tablets that a pregnant woman
received (in this case 2.0 per pregnant woman in the
population).\footnote{Surveys can provide coverage of IFA
  supplementation, as here the unit of data analysis is individual
  pregnant women.}

\begin{figure}[H]

{\centering \includegraphics[width=1\linewidth,height=\textheight,keepaspectratio]{pages/../images/2-nat-fpet.png}

}

\caption{Subnational Coverage Inequality for ANC4}

\end{figure}%

\subsection{Delivery care}\label{sec-nat-cov-delivery}

All countries have at least one delivery care indicator with a target in
the national plan. The global ENAP/EPMM coverage target for 2025 is at
least 90\% global average coverage, and 90\% of countries at least 80\%
coverage of skilled birth attendance (SBA). For postnatal care (PNC)
within 2 days the global coverage target is 80\% and the 90\% of
countries with at least 60\% coverage.

\textbf{\emph{Institutional (live) birth coverage and SBA}} are closely
related, as almost all deliveries with a skilled attendant occur in
health facilities. From the analytical point of view, institutional
birth coverage is preferred because it is a more objective measure and
avoids issues with the definition of what constitutes skilled birth. SBA
is often preferred from the health care perspective, as it includes an
element of quality: obviously, an institutional delivery without skilled
attendance is not desirable, and in some countries home SBA may be part
of the health care delivery strategy. attendance. Either indicator works
well to capture delivery care.

\textbf{\emph{Caesarean section}} is a live saving intervention and an
important indicator. A general rule of thumb is, put forward by WHO,
that in a population the need for Caesarean section is in the range of
10-15\% of all deliveries. If the Caesarean section rate is below 10\%,
that means that there is unmet need. If the Caesarean section rate is
over 10-15\%, that implies that there most likely overuse of Caesarean
section. It may however also imply that there is still unmet need among
certain population groups (e.g.~the poorest women or rural women) in
combination with overuse in other subgroups of the population (such as
urban and richer women).

\textbf{\emph{Postnatal care (PNC) within 48 hours}} is provided to the
mother/women and newborn. Country systems may have different ways of
recording the type of PNC and also surveys are known to have data issues
with PNC for the mother or the baby.\footnote{Amouzou A, Hazel E, Vaz L,
  Sanni Y, Moran A. Discordance in postnatal care between mothers and
  newborns: Measurement artifact or missed opportunity? J Glob Health.
  2020 Jun;10(1):010505.} Some countries use multiple definitions of
timing of the PNC visit (e.g., within the first week).

\textbf{\emph{Low birthweight}} is a critical indicator of neonatal
health. It is most meaningful if the distinction between prematurity and
small-for-gestational age is made, but most facility reporting systems
and most surveys do not have such data. All babies are supposed to have
been weighed right after birth and the percent of newborns weighing less
than 2500 grams is usually reported in the DHIS2 system. As general
guidance for the interpretation of the data, low birthweight prevalence
in sub-Saharan Africa was estimated at 13.9\% (95\% credibility
interval: 12.4-15.7\%) for 2020.\footnote{Okwaraji YB, et al.~National,
  regional, and global estimates of low birthweight in 2020, with trends
  from 2000: a systematic analysis. Lancet. 2024 Mar
  16;403(10431):1071-1080.}

\subsection{Immunization}\label{sec-nat-cov-immunization}

Immunization coverage indicators are included in virtually every
country's health sector monitoring plan. A general target is at least
90\% coverage for essential vaccines given in childhood and adolescence.

For the national coverage analyses, the focus is on \textbf{BCG}, first
and third doses of penta/DTP (penta1 or DTP1 and penta3 or DTP3), and
first and second doses of measles vaccine, often given in combination
with rubella vaccine. \emph{BCG and penta/DTP vaccinations} are
recommended to be given at birth (BCG), 6 weeks (penta1) and 14 weeks
(penta3). For facility data, the number of vaccinations given to infants
is used and the denominator is the number of eligible children in the
population, which is approximated as live births minus neonatal deaths.
Survey data generally provide vaccination coverage among children 12-23
months (may also include the age at which the vaccination was given -
mostly before the first birthday).\footnote{Therefore, the survey data
  on immunization roughly refer to the program performance in the year
  before the survey}

The first dose of measles vaccine is generally recommended at age 9
months. For the facility data, the recording and reporting usually
separate measles given to children under 1 year and children 1 year and
older, though the quality of recording and reporting for the age group
may vary (there may be a tendency to record measles vaccinations after
12 months as given to infants). Here, we use children who have survived
the first year of life (live births minus neonatal deaths minus
postneonatal deaths) as the denominator for measuring vaccination
coverage. The second dose of measles uses children aged 24-35 months as
denominator. This can be estimated as live births minus neonatal deaths
= 2* postneonatal deaths).

WHO and UNICEF work with countries to produce annual estimates of
immunization coverage based on all data sources. The national estimates,
called WUENIC, are published and available for 2020-2024. These time
trends are included in the analysis outputs, to compare the 2019-2024
annual estimates of the facility data produced in the workshop and the
survey results.

\subsection{Family planning}\label{sec-nat-cov-fp}

The family planning coverage estimates are derived from a collaboration
of the Countdown with Track20. Track20 has developed an advanced
estimation tool called Family Planning Estimation Tool (FPET) which
focuses on three indicators: modern contraceptive use, unmet need for
modern contraceptives and demand satisfied with modern methods. The
three indicators are closely related since demand is satisfied (this is
the true coverage indicator) = use of modern contraceptives / (unmet
need + use of modern contraceptives).

FPET uses statistical modelling that incorporates all available data
from surveys and may also use estimates obtained from facility data if
the quality is sufficient.\footnote{https://www.track20.org/pages/data\_analysis/publications/methodological/family\_planning\_estimation\_tool.php}
FPET allows for various types of survey data to be integrated into the
estimates and fits a line that pulls from the trends. This utilizes the
strength of multiple data points and minimizes the risk of comparing
different surveys.

\begin{figure}[H]

{\centering \includegraphics[width=1\linewidth,height=\textheight,keepaspectratio]{pages/../images/2-nat-cov-fp.png}

}

\caption{FP among currently married women 15-49 years}

\end{figure}%

\chapter{National Inequality}\label{national-inequality}

\section{National Inequality (MADM)}\label{sec-subnat-inequality}

\subsection{Rationale, Approach and
implementation}\label{rationale-approach-and-implementation-2}

\textbf{Rationale: Scientific basis for the analysis}

Sub-national analysis of health intervention coverage is essential for
advancing universal health coverage (UHC), which aims to ensure
equitable access to quality health services for all populations.
National averages often mask critical disparities that exist across
regions, districts, or population subgroups.

Monitoring sub-national data helps identify geographical areas where
coverage is low, signaling potential inequities in health service access
and prompting targeted interventions. This is particularly important for
drawing attention to populations who are left behind and ensuring
resources are directed where they are most needed.

\textbf{Approach: Description of analytical steps}

Here we calculate \textbf{Median Absolute Deviation from the Median
(MADM)} to quantify spread in coverage in districts within an admin1
level , compared with the coverage at the said Admin1 subnational unit.

The key statistical measures for sub-national inequalities are:

\begin{itemize}
\item
  MADM: median absolute distance from the median. This measure gives an
  indication on whether the country has been successful in reducing
  inequalities between sub-national units.
\item
  Percent of sub-national units with coverage above a specific target or
  threshold: this indicator provides information on the extent to which
  a country has been successful in reaching universal coverage at the
  sub-national level.
\end{itemize}

\textbf{\emph{Summary of district and regional performance}}

Progress towards international and national targets can be measured by
computing the percentage of regions and districts that have achieved
these targets. The goal is for all regions and districts to have met the
target. Higher percentages mean less inequality.

\textbf{Implementation: Conducting analysis in the Shiny App}

Navigate to the \textbf{National Analysis}
\textbf{=\textgreater National Inequality} and inspect the plotted
regional and district results by year, with the median absolute distance
from the median (MADM, see screenshot below), as the summary measure to
assess if inequalities have reduced.

\begin{figure}[H]

{\centering \includegraphics[width=1\linewidth,height=\textheight,keepaspectratio]{pages/../images/2-nat_inequality.png}

}

\caption{Subnational Coverage Inequality for ANC4}

\end{figure}%

\textbf{Interpretation:}

\chapter{Global coverage targets for coverage of health
services}\label{global-coverage-targets-for-coverage-of-health-services}

Global Coverage Targets''

\begin{figure}[H]

{\centering \includegraphics[width=1\linewidth,height=\textheight,keepaspectratio]{pages/../images/2-nat-cov-targets.png}

}

\caption{Subnational Coverage Inequality for ANC4}

\end{figure}%

\chapter{``Sub-national Mapping}\label{sub-national-mapping}

\section{National and sub-national mapping of health service
coverage}\label{national-and-sub-national-mapping-of-health-service-coverage}

\chapter{Equity Assessment
(Equiplots)}\label{equity-assessment-equiplots}

\section{Rationale, Approach and
implementation}\label{rationale-approach-and-implementation-3}

\textbf{Rationale: Scientific basis for the analysis}

Household surveys provide critical information on inequalities. Our
focus is on three major dimensions of inequality: area of residence,
wealth and education.

For area of residence, we focus on urban/rural areas, for wealth, we use
household wealth quintiles and for education we focus on education of
the mother. Equiplots are used to assess whether the country has made
progress since 2010 in reducing the poor rich gap or the gap between
women with no education or low education and women with higher
education.

For wealth quintiles, it can be assessed what the type of inequality is,
as all categories are of the same size:

\begin{enumerate}
\def\labelenumi{\arabic{enumi})}
\tightlist
\item
  If the richest are well ahead of all other wealth quintiles, this is
  top inequality and is sometimes also referred to as mass deprivation
  (only the rich escape)
\item
  If the coverage differences are equidistant, this is a linear pattern
  where the increasing household wealth is linearly associated with
  higher coverage
\item
  If the poorest are left behind, this is called bottom inequality, and
  means marginalization of the poorest. Each pattern requires a
  different strategy of health programming and targeting.
\end{enumerate}

\section{Interpretation of equiplots}\label{interpretation-of-equiplots}

The interpretation should focus on whether inequalities have reduced
over time and to what extent global targets for coverage have been met.
Consider your audience/s and what key messages and insights you want to
get across.

The following questions should guide and be answered by the
interpretation:

\begin{enumerate}
\def\labelenumi{\arabic{enumi})}
\tightlist
\item
  What is the level of inequality in the most recent data point?
\item
  How have inequalities changed over time?
\item
  Is there any inequality pattern that can be observed?
\item
  What will be the best approaches to reduce inequalities?
\end{enumerate}

\textbf{Implementation: Conducting analysis in the Shiny App}

This will be analysed within the \textbf{National analyses
-\textgreater Equity Assessment} section in the Shiny App

You expect to see output as below

\begin{figure}[H]

{\centering \includegraphics[width=1\linewidth,height=\textheight,keepaspectratio]{pages/../images/2-nat-equity_1-3.png}

}

\caption{Penta3 coverage Inequality by Wealth, Area and Maternal
education}

\end{figure}%

\part{Sub-national Analyses}

\chapter{Coverage}\label{coverage-1}

\chapter{Sub-national Coverage}\label{sec-subnat-coverage}

\section{Rationale, Approach and
implementation}\label{rationale-approach-and-implementation-4}

\textbf{Rationale: Scientific basis for the analysis}

Sub-national analysis of health intervention coverage is essential for
advancing universal health coverage (UHC), which aims to ensure
equitable access to quality health services for all populations.
National averages often mask critical disparities that exist across
regions, districts, or population subgroups.

Monitoring sub-national data helps identify geographical areas where
coverage is low, signaling potential inequities in health service access
and prompting targeted interventions. This is particularly important for
drawing attention to populations who are left behind and ensuring
resources are directed where they are most needed.

\textbf{Approach: Description of analytical steps}

The focus here is to assess the coverage estimates (e.g., ANC4,
institutional deliveries, Penta3) at Admin1 and Admin2 levels using the
best performing denominator for the coverage computations, as decided
from the analysis in section 2 on denominators.

There will be more ``noise'' in the sub-national data with improbably
high or low coverage rates, compared to the national analyses, and more
so in the district analyses than in the regional (admin1) analyses. This
is because district analyses are affected by small numbers (more
fluctuations which may be random or due to data quality issues) and by
the health service utilization patterns by women and children. For
instance, a municipal district may get more deliveries than an adjacent
rural district because of the location of the hospitals in the municipal
district.

\textbf{Implementation: Conducting analysis in the Shiny App}

Navigate to the \textbf{Subnational Analysis} tab
\textbf{-\textgreater Subnational Coverage} and select the admin level
of interest (Region (Admin Level1) or District. Then select the
indicator of interest from the key indicators (ANC4, Institutional
deliveries, Low Birth Weight, Penta 1 and Measles 1. You can also use
the \textbf{Custom Check} tab to select any other indicator of interest.

\begin{figure}[H]

{\centering \includegraphics[width=1\linewidth,height=\textheight,keepaspectratio]{pages/../images/3-subnat_cov.png}

}

\caption{Subnational coverage}

\end{figure}%

\section{\texorpdfstring{\textbf{Sub-national statistical summaries (One
pager)}}{Sub-national statistical summaries (One pager)}}\label{sub-national-statistical-summaries-one-pager}

The aim is to produce a one-pager for each admin1 unit (generally
region, province or county) that contains the regional information as
well as a summary of the districts within the region. We refer to this
as the regional sub-national statistical summary.

The following components are included:

\begin{itemize}
\tightlist
\item
  Summary of key demographic information for the region and the
  districts: table with expected number of births in 2024 according to
  the DHIS2 projections and according to the preferred denominators
  derived from the health facility data (based on ANC1 for live births,
  and on penta1 for immunization indicators).
\item
  Line graphs with the trend in coverage of institutional deliveries and
  penta3 vaccination that include the best estimates for 2020-2024 based
  on the health facility data, as well as the estimates from the most
  recent surveys (from 2015) for the same indicators (with confidence
  intervals if possible).
\item
  Bar graphs that compare the 2024 coverage situation in the region
  compared to all other regions, and puts the region into the lowest,
  middle or upper third of regions in terms of coverage. This is done
  for both institutional deliveries and penta3 vaccination.
\item
  Table that summarizes the coverage for institutional deliveries and
  penta3 vaccination in 2024 by district.
\item
  Other indicators can be used as prioritized by the country (e.g.,
  ANC4, measles1).
\end{itemize}

\begin{figure}[H]

{\centering \includegraphics[width=1\linewidth,height=\textheight,keepaspectratio]{pages/../images/3-subnat_one_pager.png}

}

\caption{Percentage of districts in Alaotra Mangoro}

\end{figure}%

\emph{To summarize the coverage situation according to the health
facility statistics for 2024 can be done for the regional level and
shown on a map. A composite coverage index is computed as an average in
seven mother and child health indicators: \textbf{ANC4},
\textbf{institutional live birth coverage}, \textbf{SBA}, \textbf{IPT2},
\textbf{postnatal care}, \textbf{pentavalent vaccine 3rd dosage} and
\textbf{measles 1} vaccination coverage. Equal weight is given to all
indicators.}

\chapter{Inequality}\label{inequality}

\section{Rationale, Approach and
implementation}\label{rationale-approach-and-implementation-5}

\textbf{Rationale: Scientific basis for the analysis}

Reducing geographic inequality is essential for equitable health systems
and achieving the SDGs. Subnational inequalities reveal inconsistencies
in service delivery and highlight systemic barriers to healthcare
access. Tracking inequality indicators helps assess whether health
systems are becoming more equitable and whether targeted interventions
are working.

\textbf{Approach: Description of analytical steps}

Here we calculate \textbf{Median Absolute Deviation from the Median
(MADM)} to quantify spread in coverage in districts within an admin1
level , compared with the coverage at the said Admin1 subnational unit.

The key statistical measures for sub-national inequalities are:

\begin{itemize}
\item
  MADM: median absolute distance from the median. This measure gives an
  indication on whether the country has been successful in reducing
  inequalities between sub-national units.
\item
  Percent of sub-national units with coverage above a specific target or
  threshold: this indicator provides information on the extent to which
  a country has been successful in reaching universal coverage at the
  sub-national level.
\end{itemize}

\textbf{\emph{Summary of district and regional performance}}

Progress towards international and national targets can be measured by
computing the percentage of regions and districts that have achieved
these targets. The goal is for all regions and districts to have met the
target. Higher percentages mean less inequality.

\section{\texorpdfstring{\textbf{Implementation: Conducting analysis in
the Shiny
App}}{Implementation: Conducting analysis in the Shiny App}}\label{implementation-conducting-analysis-in-the-shiny-app}

Navigate to the \textbf{Subnational Analysis} tab
\textbf{=\textgreater Subnational Inequality} and inspect the plotted
regional and district results by year, with the median absolute distance
from the median (MADM, see screenshot below), as the summary measure to
assess if inequalities have reduced.

\begin{figure}[H]

{\centering \includegraphics[width=1\linewidth,height=\textheight,keepaspectratio]{pages/../images/3-subnat_inequality.png}

}

\caption{Subnational Coverage Inequality for ANC4}

\end{figure}%

\chapter{Global coverage targets}\label{sec-subnat-cov-target}

\part{Mortality}

\chapter{Institutional Mortality}\label{institutional-mortality}

\section{Rationale, Approach and
implementation}\label{rationale-approach-and-implementation-6}

\textbf{Rationale: Scientific basis for the analysis}

Maternal and perinatal mortality in health facilities are critical
indicators of the quality of care (institutional mortality).
Institutional maternal mortality is one of the key indicators proposed
to monitor progress. Facility-based mortality statistics can also be a
critical input into the estimation of population levels of mortality.

Reporting of maternal deaths, stillbirths, and early neonatal deaths
(before discharge) may occur through the routine reporting system
(DHIS2, aggregate or individual level) or be part of a maternal and
perinatal death surveillance and response system (MPDSR).

Population mortality statistics (maternal mortality ratio, stillbirth
rates, neonatal mortality rates) rely on household surveys and censuses,
with major limitations, especially for maternal mortality and
stillbirths. A promising development is the major increase in health
facility deliveries observed in many countries.

Especially if coverage of births in health facilities is high (e.g.,
over 75\%), the facility-based statistics will become a useful input
into the estimation of population levels of mortality. The main
challenge with mortality data from health facilities is under-reporting
of deaths. Maternal deaths do not only occur during birth, but also in
pregnancy and post-partum.

Reporting of stillbirth deaths requires well maintained maternity
registers, but also cross-checks to the operation theatre registers for
cases of Caesarean section. If that is not the case, deaths may not be
recorded in the maternity register, and not reported into the national
system. In addition, maternal deaths occurring in other hospital wards
are more likely to be missed, such as deaths associated with abortion,
or post-discharge due to sepsis and other causes. Neonatal deaths after
discharge, which is often within 24 hours, are also more likely to be
missed.

A critical step is a systematic assessment of data quality issues, which
forms the basis for, if possible, adjustments, and more importantly
should guide efforts to improve reporting quality.

\textbf{Approach: Description of analytical steps}

The definitions for institutional, community and population maternal
mortality and for stillbirths are:

\begin{longtable}[]{@{}
  >{\raggedright\arraybackslash}p{(\linewidth - 4\tabcolsep) * \real{0.3333}}
  >{\raggedright\arraybackslash}p{(\linewidth - 4\tabcolsep) * \real{0.3333}}
  >{\raggedright\arraybackslash}p{(\linewidth - 4\tabcolsep) * \real{0.3333}}@{}}
\toprule\noalign{}
\begin{minipage}[b]{\linewidth}\raggedright
\textbf{Indicator}
\end{minipage} & \begin{minipage}[b]{\linewidth}\raggedright
\textbf{Numerator}
\end{minipage} & \begin{minipage}[b]{\linewidth}\raggedright
\textbf{Denominator}
\end{minipage} \\
\midrule\noalign{}
\endhead
\bottomrule\noalign{}
\endlastfoot
Institutional maternal mortality ratio (iMMR) & Number of maternal
deaths in health facilities\footnote{Maternal death is defined as any
  cause related to or aggravated by pregnancy or its management
  (excluding accidental or incidental causes) during pregnancy and
  childbirth or within 42 days of termination of pregnancy, irrespective
  of the duration and site of the pregnancy} & Number of live births in
health facilities * 100,000 \\
Population maternal mortality (MMR) & Number of maternal deaths in the
population & Number of live births in the population * 100,000 \\
Community maternal mortality ratio (comer) & Number of maternal deaths
in the community & Number of live births in the community * 100,000 \\
Institutional stillbirth rate (iSBR) & Number of stillbirths in health
facilities\footnote{A baby who dies after 28 weeks of pregnancy, but
  before or during birth, is classified as a stillbirth. Often the
  distinction is made between ante-partum (macerated) and intrapartum
  (fresh) stillbirths.} & Number of births in health facilities *
1,000 \\
Population stillbirth rate (SBR) & Number of stillbirths in the
population & Number of births in the population * 1,000 \\
Community stillbirth rate (cSBR) & Number of stillbirths in the
community & Number of births in the community * 1,000 \\
Neonatal mortality (before discharge) & Number of neonatal deaths before
discharge & Number of live births in health facilities * 1000 \\
\end{longtable}

\section{Data Analysis Components:}\label{data-analysis-components}

Data analysis has the following components:

\textbf{iMMR and iSBR review}

The annual mortality rates are computed using the unadjusted data on
reported deaths and births/live births in health facilities. We do not
adjust for reporting rates and outliers (as is done for other
interventions) because it is difficult to adjust maternal deaths and
stillbirths, where the number of deaths is small and
fluctuating.\footnote{It is however good to inspect the data and
  consider extreme outliers (more than 3 standard deviations from the
  annual, or more than 5 times the median absolute deviation -- see data
  quality section.} It is however advisable to check the data for any
extreme outliers in the annual data (e.g., numbers of deaths that
clearly indicate data entry errors) and replace these with the average
of the surrounding years.

The figure below displays the institutional maternal mortality per
100,000 live births for regions (dots) and for the country as a whole
(line and annual values) by year.

\begin{center}
\includegraphics[width=1\linewidth,height=\textheight,keepaspectratio]{pages/../images/5-mortality_mmr.png}
\end{center}

Several considerations need to be made:

\begin{itemize}
\tightlist
\item
  Are there any extreme outliers on the high side which may be due to
  major data errors which need correction?
\item
  How many regions have implausibly low iMMR which is arbitrarily
  defined as less than 25 per 100,000 live births (25 is two times the
  MMR in high-income countries of 12.5); are these more advanced regions
  where mortality is expected to be lower, or are there less-developed
  regions with low mortality, which could be an indication of major
  under reporting of deaths.
\item
  A map with the iMMR by region would be a useful addition to guide the
  interpretation of the data, especially focusing on potential data
  quality issues.
\end{itemize}

\begin{center}
\includegraphics[width=1\linewidth,height=\textheight,keepaspectratio]{pages/../images/5-mmr_inst_map.png}
\end{center}

The figure below presents the stillbirth rate per 1,000 births for
regions and the country, using the same format as for maternal
mortality. The interpretation should explain:

\begin{itemize}
\item
  How many regions have implausibly low iSBR which is defined as less
  than 6 per 1,000 births (which is two times the SBR in high income
  countries)?
\item
  Are these more advanced regions where mortality is expected to be
  lower, or is this a sign of major under-reporting of deaths in
  less-developed regions of the country.
\end{itemize}

\begin{center}
\includegraphics[width=1\linewidth,height=\textheight,keepaspectratio]{pages/../images/5-mortality_sbr.png}
\end{center}

In addition, the institutional mortality levels can be compared to the
most recent mortality estimates for the population. These population
estimates could be coming from a recent national survey or census, or we
can use the UN estimates for maternal mortality (for 2020) and
stillbirth rates (for 2021).

This is to obtain an idea of the difference between the institutional
mortality and the population mortality.~Interpretation should seek to
explain:

\begin{itemize}
\tightlist
\item
  How far is the iMMR (or iSBR) from the UN estimates of the population
  mortality, including the uncertainty range of the global estimates:
  this difference will be used further to assess the data quality.
\end{itemize}

The institutional neonatal mortality rates (per 1,000 live births) based
on reported neonatal deaths may also be graphed similar to iMMR and
iSBR, but have to be interpreted with additional caution. Almost all
babies stay at least 24 hours after delivery in the hospital but after
that many are discharged and the observation time in health facilities
is variable. Therefore, the statistic is mostly referred to as neonatal
deaths before discharge per 1,000 live births, which includes day 1,
some deaths on day 2, fewer deaths on day 3 etc.

A rough guide to assess reporting completeness is that expected
mortality of neonatal deaths before discharge in health facilities
should be at least half of neonatal mortality in the population. So for
instance, if population neonatal mortality is 20 per 1,000 live births,
we expect institutional neonatal mortality at least 10 per 1,000 live
births in the health facilities.

\begin{center}
\includegraphics[width=1\linewidth,height=\textheight,keepaspectratio]{pages/../images/5-mortality_nmr.png}
\end{center}

\section{Data Quality metrics}\label{data-quality-metrics}

\begin{itemize}
\tightlist
\item
  \textbf{\emph{Ratio stillbirth to maternal deaths in the health
  facility data at national level}}
\end{itemize}

We expect maternal mortality and stillbirth to be positively correlated
given the commonalities in causes. Based on a review of the global
estimates, historical data, and health facility studies, we expect the
ratio of stillbirth to maternal death to be in the range of 5 to 25 for
countries in sub-Saharan Africa. We compute the ratio as the number of
reported stillbirths divided by the number of reported maternal deaths
in DHIS2 or MPDSR, in a specific time period (usually a year) and raise
a ``data quality flag'' if the ratio is outside the 5-25 range.

Interpretation:

\begin{itemize}
\tightlist
\item
  If the ratio is lower than 5: under-reporting of stillbirths is likely
  greater than under-reporting of maternal deaths
\item
  If the ratio is equal or greater than 25: under-reporting of maternal
  deaths is likely to be the main issue, under-reporting of stillbirth
  less serious than for maternal deaths.
\item
  If the ratio is between 5 and 25: under-reporting of maternal deaths
  and under-reporting of stillbirths are both possible, or reporting of
  both is of good quality (this requires that the level is also in the
  expected range - component 1).
\end{itemize}

\begin{center}
\includegraphics[width=1\linewidth,height=\textheight,keepaspectratio]{pages/../images/5-mortality_ratio_md_sb.png}
\end{center}

\begin{itemize}
\tightlist
\item
  \textbf{\emph{Consistency of institutional MMR with estimated
  population MMR and community MMR}}
\end{itemize}

The completeness of reporting by health facilities can be estimated by
comparing the reported iMMR based on facility data with an expected
iMMR. The population MMR, community MMR and institutional MMR should be
consistent. There will be variation in the community to institutional
MMR between populations, but it is not likely that, for instance, the
community MMR is 1,000 when the institutional MMR is 100.

We compute an expected MMR in health facilities based on assumptions
about:

\begin{itemize}
\tightlist
\item
  MMR in the whole population (including community and institutional
  deaths): For example, the lower bound, median, upper bounds of global
  estimates for each country (2023 UN estimates) could be used, or the
  results from a recent survey.
\item
  Ratio of community to institutional maternal mortality: We use
  assumptions ranging from 1.0 (where we assume community MMR is the
  same as iMMR) to 2.0 and 3.0 (cMMR is 2 or 3 times higher than
  iMMR).\footnote{The assumptions on the range of ratios for community
    to institutional MMR were selected based on studies that measured
    both institutional and population mortality estimates (or had
    information about the percent of all maternal deaths that occurred
    at health facilities and the percent of deliveries that took place
    at health facilities.} For each country, one should consider what
  this ratio could be which may depend on the proportion of births in
  health facilities. There is some evidence that the community to
  institutional ratio increases as institutional birth rates increases,
  as well as the observed percent of births in health facilities.
\end{itemize}

The population MMR is the sum of the institutional and community MMR,
weighted by the percent of births occurring in facilities. For instance,
if the institutional MMR is 100 and the community MMR is 200, and 75\%
of births are in health facilities, then the population MMR equals
\[ 0.75 * 100 + 0.25 *200 = 125  per 100,000  live births \] This can be
expressed in the following formula:

\(M_{p} = P_{i}*M_{i} + \left( 1 - P_{i} \right)*M_{c}\)

Where;

\begin{itemize}
\tightlist
\item
  M\textsubscript{p} = maternal mortality ratio in the population;
\item
  M\textsubscript{i} = institutional maternal mortality ratio;
\item
  M\textsubscript{c}= maternal mortality ratio in the community;
\item
  P\textsubscript{i} the proportion of live births in institution
\end{itemize}

Here, we have (1) iMMR from the DHIS2 data and (2) population mortality
from the UN estimates, and can compute the community MMR as
\[(125 - 0.75 * 100) / (1-0.25) = 200 \], or

\[ M_{c} =  \frac{{(M}_{p} - {P_{i}*M}_{i})}{(1 - P_{i})}\]

In the example, the ratio community to institutional mortality
\[(Mc/Mi) \]equals 200 / 100 = 2, in other words the community mortality
is two times higher than the institutional mortality.

We can now compute the expected MMR based on 1) an estimate of
population MMR 2) the ratio Mc/Mi. For instance, if 75\% delivers in
health facilities, the population MMR is 200 and the ratio
(M\textsubscript{c}/M\textsubscript{i}) is 2, then the expected
institutional MMR is
\[ (200/(2-(2-1)*0.75) = 200 / 1.25 = 160 per 100,000 live births \]
(and the community MMR is 320).

In a formula:\footnote{An important factor affecting MMR estimates from
  health facilities, especially if the data are gathered through MPDSR,
  is that many deaths occur outside of the maternity ward (upon
  re-admission). If the DHIS2 is based on the reporting of all
  cause-specific deaths in health facilities, and the maternal death is
  correctly classified, this is not an issue.}

\begin{quote}
Expected \[Mi =Mp / (Mc/Mi -((Mc/Mi)-1) * Pi))\]
\end{quote}

Finally, the completeness of facility reporting is reported iMMR
(e.g.~based on DHIS2) divided by the expected iMMR. For instance, if the
reported MMR was 100 and the expected MMR 160, then the level of
completeness of reporting is \[100/160 * 100% = 62.5%.
\]

Below is an example of the estimated completeness of facility reporting
of maternal deaths using different scenarios. The figure shows the
results on completeness of reporting with three levels of population MMR
(lower bound, median or best estimate, upper bound) and three community
to institutional mortality ratios (0.5-2.5), shown on the X-axis.

Not all scenarios are equally relevant to each country. For instance, if
there is evidence that the population MMR is lower than the UN median,
pick the scenario with the lowest MMR (the blue line). If it is
considered that community MMR could be 1.5-2.0 times higher than
institutional MMR, then the completeness estimate is 59-64\%. The
choices are arbitrary, but it is useful to consider if the range of
estimates of completeness of facility reporting can be narrowed down by
using the most plausible scenario. As default values, the ratios of 1.0,
2.0 and 3.0 are used.

\begin{center}
\includegraphics[width=1\linewidth,height=\textheight,keepaspectratio]{pages/../images/5-mortality_mmr_completeness.png}
\end{center}

A similar approach can be used for stillbirths using all births instead
of live births. The UN global stillbirth estimates for 2021 with
uncertainty ranges can be used (lower and upper bound are 90\%
uncertainty intervals from the model). There is little research on the
community to institutional stillbirth ratio (partly because community
level stillbirth reporting is more uncertain) but it is likely that the
ratios are lower than for maternal mortality, as institutional mortality
levels are much higher for stillbirth rates than for MMR. A range of
0.5-1.5 may be used for the estimation of the level of completeness of
facility reporting.

\begin{center}
\includegraphics[width=1\linewidth,height=\textheight,keepaspectratio]{pages/../images/5-mortality_sbr_completeness.png}
\end{center}

\chapter{}\label{section5-curative-health-services-utilization-for-sick-children}

\part{Service Utilization}

\chapter{Curative health services utilization for sick
children}\label{curative-health-services-utilization-for-sick-children}

\section{Rationale, Approach and
implementation}\label{rationale-approach-and-implementation-7}

\textbf{Rationale: Scientific basis for the analysis}

There is only limited information about curative service utilization,
even though diarrhoea and pneumonia are leading causes of death in
children. Service utilization statistics on care-seeking behaviour among
children with recent illnesses (diarrhoea, acute respiratory infection,
or fever in the last 2 weeks) is usually obtained from household
surveys, relying on mother's recall.

Health facility data on outpatient (OPD) visits is an indicator of
access to curative services: \textbf{less than one visit per person is
often considered an indicator of poor access}. Similarly, data on
hospital admissions is an indicator of access to services, while
hospital mortality (case fatality) is an indicator of the quality of
care.

\textbf{Approach: Description of analytical steps}

The data on OPD visits should include new and re-visits. Data are
usually reported for under-5 years and 5 years and older. As with data
on maternal and newborn care and immunization, the data quality is
assessed, and adjustments are made for completeness of reporting and
extreme outliers are corrected. The adjusted clean data are used for
analysis at national and subnational levels.

\subsection{Outpatient service
utilization}\label{outpatient-service-utilization}

The data on outpatient visits should include both new and re-visits.

\textbf{\emph{Mean number of OPD visits per child per year:}}

\begin{itemize}
\tightlist
\item
  The numerator is the adjusted number of under-5 OPD visits in a year
  and the denominator is the total number of children under-5 which is
  taken from the DHIS2 projections. We do not expect this statistic to
  change much between years (less than 0.2 visits per child per year). A
  gradual increase suggests either improvements in access to OPD
  services or a greater disease burden for children. There is no fixed
  cut-off, but if the attendance is less than 1 visit per year per
  child, access to services is likely an issue. The OPD statistics are
  computed for the national and regional/provincial levels.
\end{itemize}

\begin{center}
\includegraphics[width=1\linewidth,height=\textheight,keepaspectratio]{pages/../images/6-utilization_opd_utilization.png}
\end{center}

\begin{itemize}
\tightlist
\item
  A map with OPD use by region or province can reveal important
  sub-national differences.
\end{itemize}

\begin{center}
\includegraphics[width=1\linewidth,height=\textheight,keepaspectratio]{pages/../images/6-opd_map.png}
\end{center}

\begin{itemize}
\item
  OPD interpretations considerations:

  \begin{itemize}
  \item
    What can be said about the data quality for OPD visits? Is there
    consistency of reported numbers between years?
  \item
    What is the number of OPD visits per child per year during
    2019-2023, is it increasing?
  \item
    Is it lower than 1 visit per year, which is considered indicative of
    low access?
  \item
    What can be said about the OPD visits per child per year by
    region/province in 2023? How large is the difference between top and
    bottom regions?
  \end{itemize}
\end{itemize}

\textbf{\emph{Percent of OPD visits that are children under five years}}

The percentage commonly lies between 15-45\% of all visits that are
under-fives. In high fertility countries (e.g., total fertility rate
\textgreater{} 4) we expect a higher percentage (e.g., over 30\%) than
in lower fertility countries. If the percent lies outside this range,
there may be a data quality problem. Furthermore, if the percent changes
much between years (e.g.~more than 5 percentage points) then there may
also be a data quality issue.

\begin{center}
\includegraphics[width=1\linewidth,height=\textheight,keepaspectratio]{pages/../images/6-utilization_perc_under5_utilization.png}
\end{center}

\subsection{Inpatient service
utilization}\label{inpatient-service-utilization}

The data on hospital admissions (or discharges + deaths) include new and
re-admissions. Data are usually reported for under-5 and 5 years and
over. Some countries report discharges rather than admissions which
would be the preferred data (discharges = admissions -- deaths).

A review of completeness of reporting and the presence of extreme
outliers is used to assess data quality. Reporting rates of hospitals
(and other facilities with in-patient services) may be more difficult to
assess than for other services. Therefore, the decision to adjust for
incomplete reporting also depends on the judgement by the country teams
regarding the quality of the reporting rate for in-patient services.

Also, extreme outliers may be more common for monthly admissions because
of poor reporting, and adjustments need to be made cautiously. It is
recommended to assess both the unadjusted and adjusted results.

\textbf{\emph{Number of admissions per 100 children under-5 per year:}}

\begin{itemize}
\item
  This is an indicator of access. A low value, e.g., less than 2
  admissions per 100 children under-5 per year, is indicative of low
  access to services. The median for countries in sub-Saharan Africa for
  2018-2022 was 4.5 admissions per year. Also here, we do not expect the
  indicator to change much per year: e.g.~a change of 1 or more
  admissions per 100 children between years is unlikely, unless a
  specific explanation can be found (such as an epidemic).
\item
\end{itemize}

\begin{center}
\includegraphics[width=1\linewidth,height=\textheight,keepaspectratio]{pages/../images/6-utilization_ipd_utilization.png}
\end{center}

\begin{itemize}
\tightlist
\item
  A map with IPD admissions among under-five use by region or province
  can reveal important sub-national differences.
\end{itemize}

\begin{center}
\includegraphics[width=1\linewidth,height=\textheight,keepaspectratio]{pages/../images/6-ipd_map.png}
\end{center}

\begin{itemize}
\item
  IPD interpretations considerations:

  \begin{itemize}
  \item
    Is there consistency of reported numbers of admissions / admission
    rates over time?
  \item
    What is the number of admissions per 100 children under 5 per year
    during 2019-2023?
  \item
    Trend - Is it low or high? What can be said about admissions per 100
    children under-5 per year by region/province in 2023?
  \end{itemize}
\end{itemize}

\textbf{\emph{Percent of admissions that are children under five:}}

This is an indicator of data quality. The percentage commonly lies
between 10-40\% of all admissions that are under-fives. If the
percentage lies outside this range, there may be a data quality problem
or an exceptional situation. Furthermore, if the percent changes much
between years (e.g.~more than 5 percentage points), then there may also
be a data quality issue.

\begin{center}
\includegraphics[width=1\linewidth,height=\textheight,keepaspectratio]{pages/../images/6-utilization_perc_under5_utilization.png}
\end{center}

\textbf{Case fatality rate:}

An indicator of the quality of care, defined as the number of children
who die in hospital divided by the total number admitted (discharges +
deaths). This should be done using the unadjusted data, as we do not
adjust: neither the numbers of deaths nor the number of admissions. The
case fatality rate is considered an indicator of the quality of care.
The lower the mortality in health facilities, the better the quality of
care.

\begin{center}
\includegraphics[width=1\linewidth,height=\textheight,keepaspectratio]{pages/../images/6-utilization_cfr_utilization.png}
\end{center}

\part{Health System Perfomance}

\chapter{Health Systems Performance}\label{health-systems-performance}

\section{Introduction}\label{introduction}

This section focuses on analyzing health systems performance, including
the availability and quality of health services, health workforce, and
health financing. The analysis aims to identify gaps and challenges in
the health system that may affect the delivery of reproductive,
maternal, newborn, child and adolescent health and nutrition (RMNCAH-N)
services.

It has the following subsections:

\begin{itemize}
\item
  Health system inputs (national and subnational)
\item
  Health system outputs (national and subnational)
\item
  Private sector and RMNCAH-N services
\end{itemize}

\section{Rationale, approach and
implementation}\label{rationale-approach-and-implementation-8}

\textbf{Rationale: Scientific basis for the analysis}

The assessment of the burden of disease, the coverage, quality and
equity of interventions and the inputs of the health system should guide
policy making and targeting of programmes. Subnational analyses are
critical: districts and regions/provinces are key units of the health
systems and their service delivery. One element of the assessment is a
basic assessment of the health system inputs in terms of financing,
health workforce and infrastructure and outputs of the system in terms
of service utilization and coverage.

\textbf{Approach: Description of analytical steps}

The focus is on a basic comparison of health system outputs (coverage of
interventions) with system inputs (infrastructure, workforce, financing)
at the sub-national level (admin1). More complex methods such as
efficiency analyses, frontier analysis, considering socioeconomic level
of development and other factors are beyond the scope of this section of
the analysis.

This section focuses on analyzing health systems performance, including
the availability and quality of health services, health workforce, and
health financing. The analysis aims to identify gaps and challenges in
the health system that may affect the delivery of reproductive,
maternal, newborn, child and adolescent health and nutrition (RMNCAH-N)
services.

It has the following subsections:

\begin{itemize}
\tightlist
\item
  Health system inputs (national and subnational)
\item
  Health system outputs (national and subnational)
\item
  Private sector and RMNCAH-N services
\end{itemize}

The analysis for these subsections is conducted within the Health System
Performance section of the Shiny App as shown here:

\begin{center}
\includegraphics[width=1\linewidth,height=\textheight,keepaspectratio]{pages/../images/6-hs_hs_performance.png}
\end{center}

\chapter{Health systems inputs}\label{health-systems-inputs}

\section{Health systems inputs}\label{health-systems-inputs-1}

First, the assessment focuses on the quality of data for the health
system indicators at national and sub-national levels. For selected
indicators, the assessment should focus on:

\begin{enumerate}
\def\labelenumi{\arabic{enumi}.}
\item
  Comparison with global data for selected indicators (national level
  only)
\item
  Plausibility of indicator values by subnational units -- major
  outliers? Improbable patterns?
\item
  In addition, it is useful to explore the associations of the health
  system indicators with each other (e.g.~workforce and beds), if only
  to detect inconsistencies by admin1 (province, region, county).
\item
  It is also useful to assess the association of the health system
  performance between different administrative levels (e.g., admin1 and
  district) to detect outliers or inconsistencies.
\end{enumerate}

\textbf{\emph{Health financing:}}

Health financing indicators at the district or admin1 levels are
difficult to obtain and are often limited to budget and not
expenditures. The data also tends to be limited to government resources
and may miss on other sources of financing. (\emph{These data are not
used here, but if available, the financial data should be used to assess
health system inputs.)}

\textbf{\emph{Core health professionals per 10,000 population:}}

Health workforce indicators are often of poor quality and not easy to
obtain. The main indicator is the number of core health professionals
per 10,000 population. These include physicians, non-physician
clinicians (depending on the country, but often with surgical skills and
multiple years of training but no academic degree), nurses and midwives.

In 2006, WHO suggested that at least 23 core health professionals were
needed to make major progress in reducing maternal and child mortality
with high skilled birth attendance. More recently, higher thresholds
have been used: at least 44.5 per 10,000 population to achieve universal
health coverage.

\begin{center}
\includegraphics[width=1\linewidth,height=\textheight,keepaspectratio]{pages/../images/7-health_system_density_hwf.png}
\end{center}

\textbf{\emph{Number of health facilities per 10,000 population:}}

Health infrastructure is another useful indicator, including all
hospitals, health centers and lower level health facilities such as
health posts and dispensaries. Both private and public sectors should be
included. The number of health facilities per 10,000 population is a
crude indicator as it mixes small and large facilities (2 per 10,000
could be used as an indicative number, where less than 2 is considered
low).

\begin{center}
\includegraphics[width=1\linewidth,height=\textheight,keepaspectratio]{pages/../images/7-health_system_density_fac.png}
\end{center}

\textbf{\emph{Number of hospitals per 100,000 population :}}

Additional insights into the infrastructure for in-patient services can
be obtained by computing the number of hospitals per 100,000 population.

\begin{center}
\includegraphics[width=1\linewidth,height=\textheight,keepaspectratio]{pages/../images/7-health_system_density_hosp.png}
\end{center}

\textbf{\emph{Inpatient bed density per 10,000 population:}}

Additional indicator for health infrastructure, computed as the number
of inpatient beds in all health facilities per 10,000 population.

\begin{center}
\includegraphics[width=1\linewidth,height=\textheight,keepaspectratio]{pages/../images/7-health_system_density_bed.png}
\end{center}

\chapter{Health systems outputs by
inputs}\label{health-systems-outputs-by-inputs}

\begin{center}\rule{0.5\linewidth}{0.5pt}\end{center}

\section{Health systems outputs by inputs at the subnational
level}\label{health-systems-outputs-by-inputs-at-the-subnational-level}

The second part of the analysis explores the association between health
system inputs and service utilization and coverage by sub-national units
(admin1).~

The following can be examined:

\begin{enumerate}
\def\labelenumi{\arabic{enumi}.}
\tightlist
\item
  \textbf{\emph{Association between hospital density (per 100,000) and
  admission rates for children under 5 and beds per 10,000 population
  and admission rates for children under 5, by admin1.}}~
\end{enumerate}

\begin{itemize}
\item
  We expect that regions with lower hospital density have lower
  admission rates for children, and those with higher density have
  higher admission rates. This would show as a positive slope of a
  linear regression line, as below.~

  \emph{Insert screenshot (6-ipd\_use)}
\item
  There may be major outliers. For instance, regions with low density
  and high admission rates: this may be because:

  \begin{itemize}
  \item
    Hospital density is under-reported by these regions
  \item
    Hospitals in the low density regions have very high admission and
    bed occupancy rates
  \item
    Hospital admission are over-reported in these regions.
  \end{itemize}

  The interpretation should be based on knowledge of the actual
  situation in the regions.
\end{itemize}

\begin{enumerate}
\def\labelenumi{\arabic{enumi}.}
\setcounter{enumi}{1}
\tightlist
\item
  \textbf{\emph{Association between health workforce (core health
  professionals per 10,000) and outpatient visits among children under
  5, by region}}
\end{enumerate}

\begin{itemize}
\tightlist
\item
  We expect that regions with lower health workforce density have lower
  OPD utilization rates for children, and those with higher density have
  higher admission rates. This would show as a positive slope of a
  linear regression line, as in the figure below.
\end{itemize}

\begin{center}
\includegraphics[width=1\linewidth,height=\textheight,keepaspectratio]{pages/../images/8-health_system_ratio_opd_u5_hwf.png}
\end{center}

\begin{enumerate}
\def\labelenumi{\arabic{enumi}.}
\item
  There may be major outliers, for instance, regions with low health
  workforce density and high OPD utilization rates. This may be because:

  \begin{itemize}
  \item
    Health workforce density is under-reported by these regions
  \item
    Health workers have a very high workload in these regions
  \item
    Health workforce density is over-reported in these regions (less
    likely).
  \end{itemize}

  The interpretation should be based on knowledge of the actual
  situation in the regions.

  \begin{enumerate}
  \def\labelenumii{\arabic{enumii}.}
  \setcounter{enumii}{2}
  \tightlist
  \item
    \textbf{\emph{Association between health workforce (core health
    professionals per 10,000) and institutional live birth coverage, by
    region}}
  \end{enumerate}
\end{enumerate}

\begin{itemize}
\item
  For delivery coverage, data by region from either a recent household
  survey or the most recent DHIS2 data derived estimate can be used.
\item
  We expect that regions with lower health workforce density have lower
  institutional delivery rates for children, and those with higher
  density have higher rates. This would show as the positive slope of a
  linear regression line.

  The considerations for the interpretation are the same as above.
\end{itemize}

\{r, out.width = ``100\%'', fig.align = ``center'', echo=F\}
knitr::include\_graphics(here::here(``images'',``8-health\_system\_ratio\_opd\_u5\_hwf.png''))

\begin{center}
\includegraphics[width=1\linewidth,height=\textheight,keepaspectratio]{pages/../images/8-health_system_instdeliveries_hwf.png}
\end{center}

\chapter{Private sector and RMNCH
services}\label{private-sector-and-rmnch-services}

\section{Private sector share
analysis}\label{private-sector-share-analysis}

\subsection{Rationale, Approach and
implementation}\label{rationale-approach-and-implementation-9}

\textbf{Rationale: Scientific basis for the analysis}

The private sector plays a significant role in delivering RMNCH
(Reproductive, Maternal, Newborn, and Child Health) services, although
its contribution varies both between and within countries. While routine
health facility data are intended to capture private sector service
delivery---such as the number of deliveries or family planning
visits---this information is often incomplete, as reporting from private
facilities tends to be less consistent than from public ones. Ideally,
all private facilities should be included in the master facility list,
but this is frequently not the case.

To estimate the private sector's contribution more accurately, survey
data such as those from DHS and MICS can be used. These surveys include
questions about the source of health services, distinguishing between
public and private providers.

\textbf{Approach: Description of analytical steps}

In this analysis, we use three key indicators to assess the private
sector's role at the national level and disaggregated by urban and rural
residence.

The three indicators are:

\begin{enumerate}
\def\labelenumi{\arabic{enumi}.}
\item
  \textbf{Deliveries}: Percent of live births occurring in health
  facilities among all live births in the last two years, by public
  private; with the share that is private
\item
  \textbf{Surgical interventions}: Percent of C-sections that occur in
  health facilities among all live births in the last two years, by
  public private; with the share that is private
\item
  \textbf{Curative care for children}: Percent of children under-five
  who have sought care for fever, acute respiratory illness or diarrhoea
  in the last two weeks before the interview, by public private; with
  the share that is private
\end{enumerate}

It is generally expected that women and children in urban areas rely
more on privately provided services than those in rural areas.

There are, however, two caveats to the interpretation:

\begin{itemize}
\item
  Rural women may use private services located in urban areas, which
  means that the percentage of women / children is overestimated for
  rural residents.
\item
  The distinction between private-for-profit and non-profit private
  facilities is not made in this analysis. Each country uses a different
  classification and naming convention, and also for the survey
  respondent this distinction is difficult to make.
\end{itemize}

\textbf{Implementation: Conducting analysis in the Shiny App}

This can be analyzed using the Health System Performance tab
-\textgreater{}

\{r, out.width = ``100\%'', fig.align = ``center'', echo=F\}
knitr::include\_graphics(here::here(``images'',``6-utilization\_perc\_under5\_utilization.png''))

\begin{center}
\includegraphics[width=1\linewidth,height=\textheight,keepaspectratio]{pages/../images/8-health_system_private_sector_general.png}
\end{center}

\bookmarksetup{startatroot}

\chapter{Planning Ahead}\label{planning-ahead}

The objective of the final session of the workshop is to identify
concrete actions to take forward the work you have completed. As you
have been preparing and interpreting your analyses, you have been
encouraged to think about the full cycle of data-driven decision-making.
(See diagram below.) This starts with identifying the data users and
understanding the questions or issues they aim to address, and then
preparing the findings, interpretation and insights that would be most
useful for addressing those questions or issues. Information products
and dissemination strategies should be tailored to align with users'
preferences and should be supplemented with support to decision-makers
in utilizing the analysis and following up on agreed actions. This
requires ongoing engagement between data providers and users to
understand needs and to learn what is most useful in responding to those
needs.

\begin{figure}[H]

{\centering \includegraphics[width=1\linewidth,height=\textheight,keepaspectratio]{pages/../images/en-7-planning_ahead.png}

}

\caption{Diagram. Cycle of data-drive decision-making}

\end{figure}%

As you prepare your plans, consider your use cases and audiences, as
well as the platforms and meetings that could be leveraged to reach
these audiences. It will be important to do more than just share the
analysis as a one-off event, so think about how to feed into routine
forums where data and analyses are used, and decisions are taken. A
primary use will be contributing to routine reviews of progress and
performance led by the Ministry of Health, such as analytical reports
for annual reviews or for global reporting on RMNCAH-N.

Consider also how this will connect with your Countdown country
collaboration and existing plans to build capacity for RMNCAH data
analysis and use. For instance, if you are having working sessions with
your target audiences to share the findings, you could also use this as
an opportunity to explore some of the concepts learned during the
workshop and to strengthen competencies to interpret the analyses and
identify opportunities for action (see further suggestions below). If
you have participated in CAM workshops in previous years, reflect on
what has worked well so far and whether there are new use cases,
audiences or dissemination approaches you could explore, drawing on your
experiences and those of other countries.

\begin{longtable}[]{@{}
  >{\raggedright\arraybackslash}p{(\linewidth - 2\tabcolsep) * \real{0.0949}}
  >{\raggedright\arraybackslash}p{(\linewidth - 2\tabcolsep) * \real{0.9051}}@{}}
\toprule\noalign{}
\begin{minipage}[b]{\linewidth}\raggedright
Purpose of the activity
\end{minipage} & \begin{minipage}[b]{\linewidth}\raggedright
Examples
\end{minipage} \\
\midrule\noalign{}
\endhead
\bottomrule\noalign{}
\endlastfoot
Supporting data use and action &
\begin{minipage}[t]{\linewidth}\raggedright
\begin{itemize}
\tightlist
\item
  Present \textbf{findings, interpretation and key insights} to the
  country platform and technical workings groups for them to review and
  interpret. Support them to identify \textbf{opportunities to take
  action based on data and evidence - for example,} address gender or
  geographic inequities, test pilots with implementation research or
  evaluation, advocate for policy change or resources, etc.
\item
  Include analysis as \textbf{part of evidence base to inform the
  development of new RMNCAH-N strategy}, a \textbf{mid-term review} and
  an \textbf{evaluation,} where any of these are underway.
\item
  Produce \textbf{an article or analytical brief for targeted audiences}
  with key learning, takeaways, and proposed actions following the
  workshop, as an entry point for follow-up.
\end{itemize}
\end{minipage} \\
Building capacity & \begin{minipage}[t]{\linewidth}\raggedright
\begin{itemize}
\tightlist
\item
  Carry out \textbf{cascade training} at national and district levels on
  the \textbf{concepts,} \textbf{interpretation, and use} of specific
  analyses, through the country Countdown collaboration.
\item
  Hold sessions to sensitize \textbf{national bureau of statistics
  staff} on Countdown methodology to build ownership.
\item
  Share the concepts and findings of \textbf{health facility data
  quality assessment} with \textbf{HMIS and M\&E staff} to build their
  understanding and as an opportunity for strengthening routine systems
  for data quality.
\item
  Build HMIS and M\&E staff capacity to understand the \textbf{relevance
  of gender and equity-related} indicators and analyses.
\item
  Build capacity of staff to incorporate \textbf{some of the new
  analytical techniques} and \textbf{approaches} into dashboards, policy
  briefs, statistical reports, etc.
\end{itemize}
\end{minipage} \\
Identifying needs for further analysis &
\begin{minipage}[t]{\linewidth}\raggedright
\begin{itemize}
\tightlist
\item
  Identify opportunities for \textbf{implementation research} to
  \textbf{address bottlenecks or to inform potential
  scale-up/replication of successes} based on findings from the
  analysis.
\item
  Identify opportunities for further research or analysis in order to
  \textbf{better understand trends in gender and inequity} that were
  observed in the analysis and to inform options for action.
\end{itemize}
\end{minipage} \\
\end{longtable}




\end{document}
